%% ============================================================
%% Lecture 5: Multiple Regression Analysis: Further Issues
%% Intermediate Econometrics — SUFE
%% Wooldridge, Introductory Econometrics, 8th edition, Chapter 6
%% ============================================================
\documentclass[aspectratio=169, 12pt]{beamer}
% ==============================================================================
% header.tex — LaTeX Preamble for Intermediate Econometrics
% Shandong University of Finance and Economics (SUFE)
% Textbook: Wooldridge, Introductory Econometrics, 8th edition
% ==============================================================================
%
% USAGE (from Slides/ directory):
%   % ==============================================================================
% header.tex — LaTeX Preamble for Intermediate Econometrics
% Shandong University of Finance and Economics (SUFE)
% Textbook: Wooldridge, Introductory Econometrics, 8th edition
% ==============================================================================
%
% USAGE (from Slides/ directory):
%   % ==============================================================================
% header.tex — LaTeX Preamble for Intermediate Econometrics
% Shandong University of Finance and Economics (SUFE)
% Textbook: Wooldridge, Introductory Econometrics, 8th edition
% ==============================================================================
%
% USAGE (from Slides/ directory):
%   \input{header}          % in Beamer .tex files (TEXINPUTS=../Preambles:)
%
% COMPILE (3-pass XeLaTeX):
%   cd Slides
%   TEXINPUTS=../Preambles:$TEXINPUTS xelatex -interaction=nonstopmode file.tex
%   BIBINPUTS=..:$BIBINPUTS bibtex file
%   TEXINPUTS=../Preambles:$TEXINPUTS xelatex -interaction=nonstopmode file.tex
%   TEXINPUTS=../Preambles:$TEXINPUTS xelatex -interaction=nonstopmode file.tex
% ==============================================================================


% --- Essential packages -------------------------------------------------------
\usepackage{fontspec}           % XeLaTeX font support (must load early)
\usepackage{amsmath}            % Math environments (align, equation, etc.)
\usepackage{amssymb}            % Math symbols (mathbb, etc.)
\usepackage{bm}                 % Bold math (\bm{})
\usepackage{mathtools}          % Extended math tools (coloneqq, etc.)


% --- Table packages -----------------------------------------------------------
\usepackage{booktabs}           % Professional tables (\toprule, \midrule, \bottomrule)
\usepackage{threeparttable}     % Table notes (tablenotes environment)
\usepackage{siunitx}            % Number alignment in tables (S column type)
\sisetup{
  table-format = 1.4,
  table-number-alignment = center,
  detect-weight = true,
  detect-inline-weight = math
}


% --- Figure and graphics packages --------------------------------------------
\usepackage{graphicx}           % Include figures (\includegraphics)
\usepackage{subcaption}         % Subfigures (subfigure environment)
\usepackage{tikz}               % TikZ diagrams
\usepackage{pgfplots}           % Data plots
\pgfplotsset{compat=1.18}
\usetikzlibrary{arrows.meta, positioning, shapes.geometric,
                decorations.pathreplacing, calc, patterns}


% --- Color and boxes ----------------------------------------------------------
\usepackage{xcolor}             % Extended color support
\usepackage{tcolorbox}          % Custom box environments
\tcbuselibrary{skins, breakable, theorems}


% --- Other utility packages ---------------------------------------------------
\usepackage{hyperref}           % Hyperlinks
\usepackage{natbib}             % Bibliography (\citet, \citep, \citealt)
\usepackage{caption}            % Figure/table captions
\usepackage{appendixnumberbeamer} % Appendix frame numbering


% ==============================================================================
% SUFE INSTITUTIONAL COLORS
% ==============================================================================
% Update hex values to match official SUFE brand colors if available.

\definecolor{SUFEBlue}{HTML}{012169}     % Primary: deep navy
\definecolor{SUFEGold}{HTML}{B9975B}     % Secondary: warm gold
\definecolor{SUFEYellow}{HTML}{F2A900}   % Highlight: bright gold/yellow
\definecolor{SUFELight}{HTML}{E8EDF5}    % Light background: pale blue-gray
\definecolor{SUFEDark}{HTML}{1a1a1a}     % Body text: near-black
\definecolor{positive}{HTML}{2E7D32}     % Positive/correct: dark green
\definecolor{negative}{HTML}{C62828}     % Negative/incorrect: dark red


% ==============================================================================
% BEAMER THEME & COLOR SETUP
% ==============================================================================

\usetheme{default}
\useinnertheme{default}
\useoutertheme{default}
\usecolortheme{default}

% Frame title
\setbeamercolor{frametitle}{fg=SUFEBlue, bg=white}
\setbeamerfont{frametitle}{size=\large, series=\bfseries}
\setbeamertemplate{frametitle}[default][left]
\addtobeamertemplate{frametitle}{\vspace{0.05em}}{%
  \vspace{0.05em}{\color{SUFEGold}\rule{\textwidth}{0.5pt}}\vspace{-0.1em}}

% Title page
\setbeamercolor{title}{fg=SUFEBlue}
\setbeamercolor{subtitle}{fg=SUFEGold}
\setbeamercolor{author}{fg=SUFEBlue}
\setbeamercolor{institute}{fg=SUFEDark}
\setbeamercolor{date}{fg=SUFEGold}

% Structure elements
\setbeamercolor{structure}{fg=SUFEBlue}
\setbeamercolor{item}{fg=SUFEBlue}
\setbeamercolor{subitem}{fg=SUFEGold}
\setbeamercolor{subsubitem}{fg=SUFEGold}
\setbeamerfont{itemize/enumerate body}{size=\normalsize}
\setbeamerfont{itemize/enumerate subbody}{size=\small}

% Block environments
\setbeamercolor{block title}{fg=white, bg=SUFEBlue}
\setbeamercolor{block body}{fg=SUFEDark, bg=SUFELight}
\setbeamercolor{block title alerted}{fg=white, bg=red!70!black}
\setbeamercolor{block body alerted}{fg=SUFEDark, bg=red!5}
\setbeamercolor{block title example}{fg=white, bg=SUFEGold!80!black}
\setbeamercolor{block body example}{fg=SUFEDark, bg=SUFEGold!8}

% Remove navigation symbols
\setbeamertemplate{navigation symbols}{}

% Footline: author | short title | slide N/N
\setbeamertemplate{footline}{%
  \leavevmode%
  \hbox{%
    \begin{beamercolorbox}[wd=.30\paperwidth,ht=2.25ex,dp=1ex,left,leftskip=2mm]{author in head/foot}%
      \usebeamerfont{author in head/foot}\insertshortauthor
    \end{beamercolorbox}%
    \begin{beamercolorbox}[wd=.40\paperwidth,ht=2.25ex,dp=1ex,center]{title in head/foot}%
      \usebeamerfont{title in head/foot}\insertshorttitle
    \end{beamercolorbox}%
    \begin{beamercolorbox}[wd=.30\paperwidth,ht=2.25ex,dp=1ex,right,rightskip=2mm]{date in head/foot}%
      \usebeamerfont{date in head/foot}%
      \insertframenumber{} / \inserttotalframenumber
    \end{beamercolorbox}}%
  \vskip0pt%
}
\setbeamercolor{author in head/foot}{fg=white, bg=SUFEBlue}
\setbeamercolor{title in head/foot}{fg=white, bg=SUFEBlue!80!black}
\setbeamercolor{date in head/foot}{fg=white, bg=SUFEBlue}

% Section separator slides (large centered section title)
\AtBeginSection[]{
  \begin{frame}
    \centering
    \vfill
    {\Large\bfseries\color{SUFEBlue}\insertsectionhead}
    \vfill
    {\color{SUFEGold}\rule{0.5\textwidth}{1pt}}
    \vfill
  \end{frame}
}


% ==============================================================================
% FONT SETUP (XeLaTeX)
% ==============================================================================
% Using Windows system fonts (Times New Roman / Arial / Consolas).
% On macOS, swap for: TeX Gyre Termes / TeX Gyre Heros / TeX Gyre Cursor
%   or: Times New Roman / Helvetica / Courier New

\setmainfont{Times New Roman}
\setsansfont{Arial}
\setmonofont{Consolas}


% ==============================================================================
% CUSTOM BOX ENVIRONMENTS (tcolorbox)
% ==============================================================================
% Quarto CSS equivalents: .keybox, .highlightbox, .methodbox,
%                         .assumptionbox, .resultbox

% keybox: key definitions, stated assumptions
\newtcolorbox{keybox}{
  enhanced,
  colback  = SUFEGold!10!white,
  colframe = SUFEGold,
  leftrule = 4pt, rightrule = 0pt, toprule = 0pt, bottomrule = 0pt,
  arc = 0pt, outer arc = 0pt,
  boxsep = 0pt, left = 8pt, right = 6pt, top = 4pt, bottom = 4pt,
}

% highlightbox: important theorems, main results
\newtcolorbox{highlightbox}{
  enhanced,
  colback  = SUFEYellow!8!white,
  colframe = SUFEYellow!80!SUFEGold,
  leftrule = 4pt, rightrule = 0pt, toprule = 0pt, bottomrule = 0pt,
  arc = 0pt, outer arc = 0pt,
  boxsep = 0pt, left = 8pt, right = 6pt, top = 4pt, bottom = 4pt,
}

% definitionbox{Title}: formal definitions — OLS, IV, GMM, etc.
\newtcolorbox{definitionbox}[1]{
  enhanced,
  title    = #1,
  colback  = SUFEBlue!4!white,
  colframe = SUFEBlue,
  colbacktitle = SUFEBlue,
  coltitle = white,
  fonttitle= \bfseries\small,
  arc = 4pt, outer arc = 4pt,
  boxsep = 0pt, left = 8pt, right = 6pt, top = 4pt, bottom = 4pt,
  toptitle = 3pt, bottomtitle = 3pt,
}

% examplebox{Title}: empirical examples, Wooldridge data applications
\newtcolorbox{examplebox}[1]{
  enhanced,
  title    = #1,
  colback  = SUFELight,
  colframe = SUFEGold!70!white,
  colbacktitle = SUFEGold!35!white,
  coltitle = SUFEBlue,
  fonttitle= \bfseries\small,
  arc = 4pt, outer arc = 4pt,
  boxsep = 0pt, left = 8pt, right = 6pt, top = 4pt, bottom = 4pt,
  toptitle = 3pt, bottomtitle = 3pt,
}

% assumptionbox: numbered OLS assumptions (MLR.1–MLR.6)
\newtcolorbox{assumptionbox}{
  enhanced,
  colback  = SUFEGold!8!white,
  colframe = SUFEGold,
  arc = 4pt, outer arc = 4pt,
  boxsep = 0pt, left = 8pt, right = 6pt, top = 4pt, bottom = 4pt,
}


% ==============================================================================
% STANDARD MATH MACROS
% ==============================================================================

% --- Operators ----------------------------------------------------------------
\DeclareMathOperator{\E}{\mathbb{E}}       % Expectation
\DeclareMathOperator{\Var}{Var}             % Variance
\DeclareMathOperator{\Cov}{Cov}             % Covariance
\DeclareMathOperator{\Corr}{Corr}           % Correlation
\DeclareMathOperator{\plim}{plim}           % Probability limit
\DeclareMathOperator{\rank}{rank}           % Matrix rank
\DeclareMathOperator{\tr}{tr}               % Matrix trace
\DeclareMathOperator{\se}{se}               % Standard error

% --- Common shortcuts ---------------------------------------------------------
\newcommand{\bhat}{\hat{\beta}}             % OLS point estimator
\newcommand{\bhatt}{\tilde{\beta}}          % Alternative estimator
\newcommand{\eps}{\varepsilon}              % Error term
\newcommand{\epsh}{\hat{\varepsilon}}       % OLS residual
\newcommand{\uh}{\hat{u}}                   % Residual (alternative)
\newcommand{\yh}{\hat{y}}                   % Fitted value

% Matrices and vectors
\newcommand{\Xb}{\mathbf{X}}               % Design matrix
\newcommand{\yb}{\mathbf{y}}               % Outcome vector
\newcommand{\betab}{\boldsymbol{\beta}}    % Parameter vector
\newcommand{\epsb}{\boldsymbol{\varepsilon}} % Error vector

% Goodness-of-fit
\newcommand{\SSR}{\text{SSR}}              % Sum of squared residuals
\newcommand{\SSE}{\text{SSE}}              % Explained sum of squares
\newcommand{\SST}{\text{SST}}              % Total sum of squares
\newcommand{\Rsq}{R^2}                     % R-squared
\newcommand{\aRsq}{\bar{R}^2}             % Adjusted R-squared

% Asymptotic notation
\newcommand{\avar}{\text{Avar}}            % Asymptotic variance
\newcommand{\pto}{\overset{p}{\to}}        % Convergence in probability
\newcommand{\dto}{\overset{d}{\to}}        % Convergence in distribution
\newcommand{\iid}{\overset{\text{iid}}{\sim}}  % i.i.d. distributed
\newcommand{\asim}{\overset{a}{\sim}}      % Approximately distributed

% Econometric estimators
\newcommand{\OLS}{\text{OLS}}
\newcommand{\IV}{\text{IV}}
\newcommand{\TSLS}{\text{2SLS}}
\newcommand{\GMM}{\text{GMM}}
\newcommand{\FE}{\text{FE}}
\newcommand{\RE}{\text{RE}}
\newcommand{\DID}{\text{DiD}}
\newcommand{\RD}{\text{RD}}

% Probability
\newcommand{\prob}{\mathbb{P}}             % Probability
\newcommand{\indep}{\perp\!\!\!\perp}      % Statistical independence

% Distributions
\newcommand{\Normal}{\mathcal{N}}          % Normal distribution
% Usage: X \sim \Normal(\mu, \sigma^2)

% Absolute value and norm
\newcommand{\abs}[1]{\left\lvert #1 \right\rvert}
\newcommand{\norm}[1]{\left\lVert #1 \right\rVert}


% ==============================================================================
% HYPERREF SETUP
% ==============================================================================

\hypersetup{
  colorlinks = true,
  linkcolor  = SUFEBlue,
  citecolor  = SUFEGold,
  urlcolor   = SUFEBlue
}


% ==============================================================================
% BIBLIOGRAPHY SETUP
% ==============================================================================

\bibliographystyle{aer}   % American Economic Review style
% Usage: \bibliography{../Bibliography_base}  (from Slides/ directory)
% In-text: \citet{Wooldridge2020}, \citep{White1980}
          % in Beamer .tex files (TEXINPUTS=../Preambles:)
%
% COMPILE (3-pass XeLaTeX):
%   cd Slides
%   TEXINPUTS=../Preambles:$TEXINPUTS xelatex -interaction=nonstopmode file.tex
%   BIBINPUTS=..:$BIBINPUTS bibtex file
%   TEXINPUTS=../Preambles:$TEXINPUTS xelatex -interaction=nonstopmode file.tex
%   TEXINPUTS=../Preambles:$TEXINPUTS xelatex -interaction=nonstopmode file.tex
% ==============================================================================


% --- Essential packages -------------------------------------------------------
\usepackage{fontspec}           % XeLaTeX font support (must load early)
\usepackage{amsmath}            % Math environments (align, equation, etc.)
\usepackage{amssymb}            % Math symbols (mathbb, etc.)
\usepackage{bm}                 % Bold math (\bm{})
\usepackage{mathtools}          % Extended math tools (coloneqq, etc.)


% --- Table packages -----------------------------------------------------------
\usepackage{booktabs}           % Professional tables (\toprule, \midrule, \bottomrule)
\usepackage{threeparttable}     % Table notes (tablenotes environment)
\usepackage{siunitx}            % Number alignment in tables (S column type)
\sisetup{
  table-format = 1.4,
  table-number-alignment = center,
  detect-weight = true,
  detect-inline-weight = math
}


% --- Figure and graphics packages --------------------------------------------
\usepackage{graphicx}           % Include figures (\includegraphics)
\usepackage{subcaption}         % Subfigures (subfigure environment)
\usepackage{tikz}               % TikZ diagrams
\usepackage{pgfplots}           % Data plots
\pgfplotsset{compat=1.18}
\usetikzlibrary{arrows.meta, positioning, shapes.geometric,
                decorations.pathreplacing, calc, patterns}


% --- Color and boxes ----------------------------------------------------------
\usepackage{xcolor}             % Extended color support
\usepackage{tcolorbox}          % Custom box environments
\tcbuselibrary{skins, breakable, theorems}


% --- Other utility packages ---------------------------------------------------
\usepackage{hyperref}           % Hyperlinks
\usepackage{natbib}             % Bibliography (\citet, \citep, \citealt)
\usepackage{caption}            % Figure/table captions
\usepackage{appendixnumberbeamer} % Appendix frame numbering


% ==============================================================================
% SUFE INSTITUTIONAL COLORS
% ==============================================================================
% Update hex values to match official SUFE brand colors if available.

\definecolor{SUFEBlue}{HTML}{012169}     % Primary: deep navy
\definecolor{SUFEGold}{HTML}{B9975B}     % Secondary: warm gold
\definecolor{SUFEYellow}{HTML}{F2A900}   % Highlight: bright gold/yellow
\definecolor{SUFELight}{HTML}{E8EDF5}    % Light background: pale blue-gray
\definecolor{SUFEDark}{HTML}{1a1a1a}     % Body text: near-black
\definecolor{positive}{HTML}{2E7D32}     % Positive/correct: dark green
\definecolor{negative}{HTML}{C62828}     % Negative/incorrect: dark red


% ==============================================================================
% BEAMER THEME & COLOR SETUP
% ==============================================================================

\usetheme{default}
\useinnertheme{default}
\useoutertheme{default}
\usecolortheme{default}

% Frame title
\setbeamercolor{frametitle}{fg=SUFEBlue, bg=white}
\setbeamerfont{frametitle}{size=\large, series=\bfseries}
\setbeamertemplate{frametitle}[default][left]
\addtobeamertemplate{frametitle}{\vspace{0.05em}}{%
  \vspace{0.05em}{\color{SUFEGold}\rule{\textwidth}{0.5pt}}\vspace{-0.1em}}

% Title page
\setbeamercolor{title}{fg=SUFEBlue}
\setbeamercolor{subtitle}{fg=SUFEGold}
\setbeamercolor{author}{fg=SUFEBlue}
\setbeamercolor{institute}{fg=SUFEDark}
\setbeamercolor{date}{fg=SUFEGold}

% Structure elements
\setbeamercolor{structure}{fg=SUFEBlue}
\setbeamercolor{item}{fg=SUFEBlue}
\setbeamercolor{subitem}{fg=SUFEGold}
\setbeamercolor{subsubitem}{fg=SUFEGold}
\setbeamerfont{itemize/enumerate body}{size=\normalsize}
\setbeamerfont{itemize/enumerate subbody}{size=\small}

% Block environments
\setbeamercolor{block title}{fg=white, bg=SUFEBlue}
\setbeamercolor{block body}{fg=SUFEDark, bg=SUFELight}
\setbeamercolor{block title alerted}{fg=white, bg=red!70!black}
\setbeamercolor{block body alerted}{fg=SUFEDark, bg=red!5}
\setbeamercolor{block title example}{fg=white, bg=SUFEGold!80!black}
\setbeamercolor{block body example}{fg=SUFEDark, bg=SUFEGold!8}

% Remove navigation symbols
\setbeamertemplate{navigation symbols}{}

% Footline: author | short title | slide N/N
\setbeamertemplate{footline}{%
  \leavevmode%
  \hbox{%
    \begin{beamercolorbox}[wd=.30\paperwidth,ht=2.25ex,dp=1ex,left,leftskip=2mm]{author in head/foot}%
      \usebeamerfont{author in head/foot}\insertshortauthor
    \end{beamercolorbox}%
    \begin{beamercolorbox}[wd=.40\paperwidth,ht=2.25ex,dp=1ex,center]{title in head/foot}%
      \usebeamerfont{title in head/foot}\insertshorttitle
    \end{beamercolorbox}%
    \begin{beamercolorbox}[wd=.30\paperwidth,ht=2.25ex,dp=1ex,right,rightskip=2mm]{date in head/foot}%
      \usebeamerfont{date in head/foot}%
      \insertframenumber{} / \inserttotalframenumber
    \end{beamercolorbox}}%
  \vskip0pt%
}
\setbeamercolor{author in head/foot}{fg=white, bg=SUFEBlue}
\setbeamercolor{title in head/foot}{fg=white, bg=SUFEBlue!80!black}
\setbeamercolor{date in head/foot}{fg=white, bg=SUFEBlue}

% Section separator slides (large centered section title)
\AtBeginSection[]{
  \begin{frame}
    \centering
    \vfill
    {\Large\bfseries\color{SUFEBlue}\insertsectionhead}
    \vfill
    {\color{SUFEGold}\rule{0.5\textwidth}{1pt}}
    \vfill
  \end{frame}
}


% ==============================================================================
% FONT SETUP (XeLaTeX)
% ==============================================================================
% Using Windows system fonts (Times New Roman / Arial / Consolas).
% On macOS, swap for: TeX Gyre Termes / TeX Gyre Heros / TeX Gyre Cursor
%   or: Times New Roman / Helvetica / Courier New

\setmainfont{Times New Roman}
\setsansfont{Arial}
\setmonofont{Consolas}


% ==============================================================================
% CUSTOM BOX ENVIRONMENTS (tcolorbox)
% ==============================================================================
% Quarto CSS equivalents: .keybox, .highlightbox, .methodbox,
%                         .assumptionbox, .resultbox

% keybox: key definitions, stated assumptions
\newtcolorbox{keybox}{
  enhanced,
  colback  = SUFEGold!10!white,
  colframe = SUFEGold,
  leftrule = 4pt, rightrule = 0pt, toprule = 0pt, bottomrule = 0pt,
  arc = 0pt, outer arc = 0pt,
  boxsep = 0pt, left = 8pt, right = 6pt, top = 4pt, bottom = 4pt,
}

% highlightbox: important theorems, main results
\newtcolorbox{highlightbox}{
  enhanced,
  colback  = SUFEYellow!8!white,
  colframe = SUFEYellow!80!SUFEGold,
  leftrule = 4pt, rightrule = 0pt, toprule = 0pt, bottomrule = 0pt,
  arc = 0pt, outer arc = 0pt,
  boxsep = 0pt, left = 8pt, right = 6pt, top = 4pt, bottom = 4pt,
}

% definitionbox{Title}: formal definitions — OLS, IV, GMM, etc.
\newtcolorbox{definitionbox}[1]{
  enhanced,
  title    = #1,
  colback  = SUFEBlue!4!white,
  colframe = SUFEBlue,
  colbacktitle = SUFEBlue,
  coltitle = white,
  fonttitle= \bfseries\small,
  arc = 4pt, outer arc = 4pt,
  boxsep = 0pt, left = 8pt, right = 6pt, top = 4pt, bottom = 4pt,
  toptitle = 3pt, bottomtitle = 3pt,
}

% examplebox{Title}: empirical examples, Wooldridge data applications
\newtcolorbox{examplebox}[1]{
  enhanced,
  title    = #1,
  colback  = SUFELight,
  colframe = SUFEGold!70!white,
  colbacktitle = SUFEGold!35!white,
  coltitle = SUFEBlue,
  fonttitle= \bfseries\small,
  arc = 4pt, outer arc = 4pt,
  boxsep = 0pt, left = 8pt, right = 6pt, top = 4pt, bottom = 4pt,
  toptitle = 3pt, bottomtitle = 3pt,
}

% assumptionbox: numbered OLS assumptions (MLR.1–MLR.6)
\newtcolorbox{assumptionbox}{
  enhanced,
  colback  = SUFEGold!8!white,
  colframe = SUFEGold,
  arc = 4pt, outer arc = 4pt,
  boxsep = 0pt, left = 8pt, right = 6pt, top = 4pt, bottom = 4pt,
}


% ==============================================================================
% STANDARD MATH MACROS
% ==============================================================================

% --- Operators ----------------------------------------------------------------
\DeclareMathOperator{\E}{\mathbb{E}}       % Expectation
\DeclareMathOperator{\Var}{Var}             % Variance
\DeclareMathOperator{\Cov}{Cov}             % Covariance
\DeclareMathOperator{\Corr}{Corr}           % Correlation
\DeclareMathOperator{\plim}{plim}           % Probability limit
\DeclareMathOperator{\rank}{rank}           % Matrix rank
\DeclareMathOperator{\tr}{tr}               % Matrix trace
\DeclareMathOperator{\se}{se}               % Standard error

% --- Common shortcuts ---------------------------------------------------------
\newcommand{\bhat}{\hat{\beta}}             % OLS point estimator
\newcommand{\bhatt}{\tilde{\beta}}          % Alternative estimator
\newcommand{\eps}{\varepsilon}              % Error term
\newcommand{\epsh}{\hat{\varepsilon}}       % OLS residual
\newcommand{\uh}{\hat{u}}                   % Residual (alternative)
\newcommand{\yh}{\hat{y}}                   % Fitted value

% Matrices and vectors
\newcommand{\Xb}{\mathbf{X}}               % Design matrix
\newcommand{\yb}{\mathbf{y}}               % Outcome vector
\newcommand{\betab}{\boldsymbol{\beta}}    % Parameter vector
\newcommand{\epsb}{\boldsymbol{\varepsilon}} % Error vector

% Goodness-of-fit
\newcommand{\SSR}{\text{SSR}}              % Sum of squared residuals
\newcommand{\SSE}{\text{SSE}}              % Explained sum of squares
\newcommand{\SST}{\text{SST}}              % Total sum of squares
\newcommand{\Rsq}{R^2}                     % R-squared
\newcommand{\aRsq}{\bar{R}^2}             % Adjusted R-squared

% Asymptotic notation
\newcommand{\avar}{\text{Avar}}            % Asymptotic variance
\newcommand{\pto}{\overset{p}{\to}}        % Convergence in probability
\newcommand{\dto}{\overset{d}{\to}}        % Convergence in distribution
\newcommand{\iid}{\overset{\text{iid}}{\sim}}  % i.i.d. distributed
\newcommand{\asim}{\overset{a}{\sim}}      % Approximately distributed

% Econometric estimators
\newcommand{\OLS}{\text{OLS}}
\newcommand{\IV}{\text{IV}}
\newcommand{\TSLS}{\text{2SLS}}
\newcommand{\GMM}{\text{GMM}}
\newcommand{\FE}{\text{FE}}
\newcommand{\RE}{\text{RE}}
\newcommand{\DID}{\text{DiD}}
\newcommand{\RD}{\text{RD}}

% Probability
\newcommand{\prob}{\mathbb{P}}             % Probability
\newcommand{\indep}{\perp\!\!\!\perp}      % Statistical independence

% Distributions
\newcommand{\Normal}{\mathcal{N}}          % Normal distribution
% Usage: X \sim \Normal(\mu, \sigma^2)

% Absolute value and norm
\newcommand{\abs}[1]{\left\lvert #1 \right\rvert}
\newcommand{\norm}[1]{\left\lVert #1 \right\rVert}


% ==============================================================================
% HYPERREF SETUP
% ==============================================================================

\hypersetup{
  colorlinks = true,
  linkcolor  = SUFEBlue,
  citecolor  = SUFEGold,
  urlcolor   = SUFEBlue
}


% ==============================================================================
% BIBLIOGRAPHY SETUP
% ==============================================================================

\bibliographystyle{aer}   % American Economic Review style
% Usage: \bibliography{../Bibliography_base}  (from Slides/ directory)
% In-text: \citet{Wooldridge2020}, \citep{White1980}
          % in Beamer .tex files (TEXINPUTS=../Preambles:)
%
% COMPILE (3-pass XeLaTeX):
%   cd Slides
%   TEXINPUTS=../Preambles:$TEXINPUTS xelatex -interaction=nonstopmode file.tex
%   BIBINPUTS=..:$BIBINPUTS bibtex file
%   TEXINPUTS=../Preambles:$TEXINPUTS xelatex -interaction=nonstopmode file.tex
%   TEXINPUTS=../Preambles:$TEXINPUTS xelatex -interaction=nonstopmode file.tex
% ==============================================================================


% --- Essential packages -------------------------------------------------------
\usepackage{fontspec}           % XeLaTeX font support (must load early)
\usepackage{amsmath}            % Math environments (align, equation, etc.)
\usepackage{amssymb}            % Math symbols (mathbb, etc.)
\usepackage{bm}                 % Bold math (\bm{})
\usepackage{mathtools}          % Extended math tools (coloneqq, etc.)


% --- Table packages -----------------------------------------------------------
\usepackage{booktabs}           % Professional tables (\toprule, \midrule, \bottomrule)
\usepackage{threeparttable}     % Table notes (tablenotes environment)
\usepackage{siunitx}            % Number alignment in tables (S column type)
\sisetup{
  table-format = 1.4,
  table-number-alignment = center,
  detect-weight = true,
  detect-inline-weight = math
}


% --- Figure and graphics packages --------------------------------------------
\usepackage{graphicx}           % Include figures (\includegraphics)
\usepackage{subcaption}         % Subfigures (subfigure environment)
\usepackage{tikz}               % TikZ diagrams
\usepackage{pgfplots}           % Data plots
\pgfplotsset{compat=1.18}
\usetikzlibrary{arrows.meta, positioning, shapes.geometric,
                decorations.pathreplacing, calc, patterns}


% --- Color and boxes ----------------------------------------------------------
\usepackage{xcolor}             % Extended color support
\usepackage{tcolorbox}          % Custom box environments
\tcbuselibrary{skins, breakable, theorems}


% --- Other utility packages ---------------------------------------------------
\usepackage{hyperref}           % Hyperlinks
\usepackage{natbib}             % Bibliography (\citet, \citep, \citealt)
\usepackage{caption}            % Figure/table captions
\usepackage{appendixnumberbeamer} % Appendix frame numbering


% ==============================================================================
% SUFE INSTITUTIONAL COLORS
% ==============================================================================
% Update hex values to match official SUFE brand colors if available.

\definecolor{SUFEBlue}{HTML}{012169}     % Primary: deep navy
\definecolor{SUFEGold}{HTML}{B9975B}     % Secondary: warm gold
\definecolor{SUFEYellow}{HTML}{F2A900}   % Highlight: bright gold/yellow
\definecolor{SUFELight}{HTML}{E8EDF5}    % Light background: pale blue-gray
\definecolor{SUFEDark}{HTML}{1a1a1a}     % Body text: near-black
\definecolor{positive}{HTML}{2E7D32}     % Positive/correct: dark green
\definecolor{negative}{HTML}{C62828}     % Negative/incorrect: dark red


% ==============================================================================
% BEAMER THEME & COLOR SETUP
% ==============================================================================

\usetheme{default}
\useinnertheme{default}
\useoutertheme{default}
\usecolortheme{default}

% Frame title
\setbeamercolor{frametitle}{fg=SUFEBlue, bg=white}
\setbeamerfont{frametitle}{size=\large, series=\bfseries}
\setbeamertemplate{frametitle}[default][left]
\addtobeamertemplate{frametitle}{\vspace{0.05em}}{%
  \vspace{0.05em}{\color{SUFEGold}\rule{\textwidth}{0.5pt}}\vspace{-0.1em}}

% Title page
\setbeamercolor{title}{fg=SUFEBlue}
\setbeamercolor{subtitle}{fg=SUFEGold}
\setbeamercolor{author}{fg=SUFEBlue}
\setbeamercolor{institute}{fg=SUFEDark}
\setbeamercolor{date}{fg=SUFEGold}

% Structure elements
\setbeamercolor{structure}{fg=SUFEBlue}
\setbeamercolor{item}{fg=SUFEBlue}
\setbeamercolor{subitem}{fg=SUFEGold}
\setbeamercolor{subsubitem}{fg=SUFEGold}
\setbeamerfont{itemize/enumerate body}{size=\normalsize}
\setbeamerfont{itemize/enumerate subbody}{size=\small}

% Block environments
\setbeamercolor{block title}{fg=white, bg=SUFEBlue}
\setbeamercolor{block body}{fg=SUFEDark, bg=SUFELight}
\setbeamercolor{block title alerted}{fg=white, bg=red!70!black}
\setbeamercolor{block body alerted}{fg=SUFEDark, bg=red!5}
\setbeamercolor{block title example}{fg=white, bg=SUFEGold!80!black}
\setbeamercolor{block body example}{fg=SUFEDark, bg=SUFEGold!8}

% Remove navigation symbols
\setbeamertemplate{navigation symbols}{}

% Footline: author | short title | slide N/N
\setbeamertemplate{footline}{%
  \leavevmode%
  \hbox{%
    \begin{beamercolorbox}[wd=.30\paperwidth,ht=2.25ex,dp=1ex,left,leftskip=2mm]{author in head/foot}%
      \usebeamerfont{author in head/foot}\insertshortauthor
    \end{beamercolorbox}%
    \begin{beamercolorbox}[wd=.40\paperwidth,ht=2.25ex,dp=1ex,center]{title in head/foot}%
      \usebeamerfont{title in head/foot}\insertshorttitle
    \end{beamercolorbox}%
    \begin{beamercolorbox}[wd=.30\paperwidth,ht=2.25ex,dp=1ex,right,rightskip=2mm]{date in head/foot}%
      \usebeamerfont{date in head/foot}%
      \insertframenumber{} / \inserttotalframenumber
    \end{beamercolorbox}}%
  \vskip0pt%
}
\setbeamercolor{author in head/foot}{fg=white, bg=SUFEBlue}
\setbeamercolor{title in head/foot}{fg=white, bg=SUFEBlue!80!black}
\setbeamercolor{date in head/foot}{fg=white, bg=SUFEBlue}

% Section separator slides (large centered section title)
\AtBeginSection[]{
  \begin{frame}
    \centering
    \vfill
    {\Large\bfseries\color{SUFEBlue}\insertsectionhead}
    \vfill
    {\color{SUFEGold}\rule{0.5\textwidth}{1pt}}
    \vfill
  \end{frame}
}


% ==============================================================================
% FONT SETUP (XeLaTeX)
% ==============================================================================
% Using Windows system fonts (Times New Roman / Arial / Consolas).
% On macOS, swap for: TeX Gyre Termes / TeX Gyre Heros / TeX Gyre Cursor
%   or: Times New Roman / Helvetica / Courier New

\setmainfont{Times New Roman}
\setsansfont{Arial}
\setmonofont{Consolas}


% ==============================================================================
% CUSTOM BOX ENVIRONMENTS (tcolorbox)
% ==============================================================================
% Quarto CSS equivalents: .keybox, .highlightbox, .methodbox,
%                         .assumptionbox, .resultbox

% keybox: key definitions, stated assumptions
\newtcolorbox{keybox}{
  enhanced,
  colback  = SUFEGold!10!white,
  colframe = SUFEGold,
  leftrule = 4pt, rightrule = 0pt, toprule = 0pt, bottomrule = 0pt,
  arc = 0pt, outer arc = 0pt,
  boxsep = 0pt, left = 8pt, right = 6pt, top = 4pt, bottom = 4pt,
}

% highlightbox: important theorems, main results
\newtcolorbox{highlightbox}{
  enhanced,
  colback  = SUFEYellow!8!white,
  colframe = SUFEYellow!80!SUFEGold,
  leftrule = 4pt, rightrule = 0pt, toprule = 0pt, bottomrule = 0pt,
  arc = 0pt, outer arc = 0pt,
  boxsep = 0pt, left = 8pt, right = 6pt, top = 4pt, bottom = 4pt,
}

% definitionbox{Title}: formal definitions — OLS, IV, GMM, etc.
\newtcolorbox{definitionbox}[1]{
  enhanced,
  title    = #1,
  colback  = SUFEBlue!4!white,
  colframe = SUFEBlue,
  colbacktitle = SUFEBlue,
  coltitle = white,
  fonttitle= \bfseries\small,
  arc = 4pt, outer arc = 4pt,
  boxsep = 0pt, left = 8pt, right = 6pt, top = 4pt, bottom = 4pt,
  toptitle = 3pt, bottomtitle = 3pt,
}

% examplebox{Title}: empirical examples, Wooldridge data applications
\newtcolorbox{examplebox}[1]{
  enhanced,
  title    = #1,
  colback  = SUFELight,
  colframe = SUFEGold!70!white,
  colbacktitle = SUFEGold!35!white,
  coltitle = SUFEBlue,
  fonttitle= \bfseries\small,
  arc = 4pt, outer arc = 4pt,
  boxsep = 0pt, left = 8pt, right = 6pt, top = 4pt, bottom = 4pt,
  toptitle = 3pt, bottomtitle = 3pt,
}

% assumptionbox: numbered OLS assumptions (MLR.1–MLR.6)
\newtcolorbox{assumptionbox}{
  enhanced,
  colback  = SUFEGold!8!white,
  colframe = SUFEGold,
  arc = 4pt, outer arc = 4pt,
  boxsep = 0pt, left = 8pt, right = 6pt, top = 4pt, bottom = 4pt,
}


% ==============================================================================
% STANDARD MATH MACROS
% ==============================================================================

% --- Operators ----------------------------------------------------------------
\DeclareMathOperator{\E}{\mathbb{E}}       % Expectation
\DeclareMathOperator{\Var}{Var}             % Variance
\DeclareMathOperator{\Cov}{Cov}             % Covariance
\DeclareMathOperator{\Corr}{Corr}           % Correlation
\DeclareMathOperator{\plim}{plim}           % Probability limit
\DeclareMathOperator{\rank}{rank}           % Matrix rank
\DeclareMathOperator{\tr}{tr}               % Matrix trace
\DeclareMathOperator{\se}{se}               % Standard error

% --- Common shortcuts ---------------------------------------------------------
\newcommand{\bhat}{\hat{\beta}}             % OLS point estimator
\newcommand{\bhatt}{\tilde{\beta}}          % Alternative estimator
\newcommand{\eps}{\varepsilon}              % Error term
\newcommand{\epsh}{\hat{\varepsilon}}       % OLS residual
\newcommand{\uh}{\hat{u}}                   % Residual (alternative)
\newcommand{\yh}{\hat{y}}                   % Fitted value

% Matrices and vectors
\newcommand{\Xb}{\mathbf{X}}               % Design matrix
\newcommand{\yb}{\mathbf{y}}               % Outcome vector
\newcommand{\betab}{\boldsymbol{\beta}}    % Parameter vector
\newcommand{\epsb}{\boldsymbol{\varepsilon}} % Error vector

% Goodness-of-fit
\newcommand{\SSR}{\text{SSR}}              % Sum of squared residuals
\newcommand{\SSE}{\text{SSE}}              % Explained sum of squares
\newcommand{\SST}{\text{SST}}              % Total sum of squares
\newcommand{\Rsq}{R^2}                     % R-squared
\newcommand{\aRsq}{\bar{R}^2}             % Adjusted R-squared

% Asymptotic notation
\newcommand{\avar}{\text{Avar}}            % Asymptotic variance
\newcommand{\pto}{\overset{p}{\to}}        % Convergence in probability
\newcommand{\dto}{\overset{d}{\to}}        % Convergence in distribution
\newcommand{\iid}{\overset{\text{iid}}{\sim}}  % i.i.d. distributed
\newcommand{\asim}{\overset{a}{\sim}}      % Approximately distributed

% Econometric estimators
\newcommand{\OLS}{\text{OLS}}
\newcommand{\IV}{\text{IV}}
\newcommand{\TSLS}{\text{2SLS}}
\newcommand{\GMM}{\text{GMM}}
\newcommand{\FE}{\text{FE}}
\newcommand{\RE}{\text{RE}}
\newcommand{\DID}{\text{DiD}}
\newcommand{\RD}{\text{RD}}

% Probability
\newcommand{\prob}{\mathbb{P}}             % Probability
\newcommand{\indep}{\perp\!\!\!\perp}      % Statistical independence

% Distributions
\newcommand{\Normal}{\mathcal{N}}          % Normal distribution
% Usage: X \sim \Normal(\mu, \sigma^2)

% Absolute value and norm
\newcommand{\abs}[1]{\left\lvert #1 \right\rvert}
\newcommand{\norm}[1]{\left\lVert #1 \right\rVert}


% ==============================================================================
% HYPERREF SETUP
% ==============================================================================

\hypersetup{
  colorlinks = true,
  linkcolor  = SUFEBlue,
  citecolor  = SUFEGold,
  urlcolor   = SUFEBlue
}


% ==============================================================================
% BIBLIOGRAPHY SETUP
% ==============================================================================

\bibliographystyle{aer}   % American Economic Review style
% Usage: \bibliography{../Bibliography_base}  (from Slides/ directory)
% In-text: \citet{Wooldridge2020}, \citep{White1980}


\title[L5: Further Issues]{Multiple Regression Analysis: Further Issues}
\subtitle{Intermediate Econometrics}
\author[SUFE Econometrics]{Chengye Jia}
\institute{%
  Department of Economics \\
  Shandong University of Finance and Economics}
\date{Spring 2026 \\ \smallskip \scriptsize Wooldridge, Chapter 6}

%% ============================================================
\begin{document}
%% ============================================================

\begin{frame}[plain]
  \titlepage
\end{frame}

\begin{frame}[shrink=10]{Outline}
  \tableofcontents
\end{frame}

%% ============================================================
\section{Data Scaling and Beta Coefficients (\S6-1)}
%% ============================================================

\begin{frame}{Data Scaling and Beta Coefficients: Motivation \& Setup}

  \begin{highlightbox}
    \textbf{Motivation.} Does it matter whether wage is in dollars or thousands, or
    birth weight in ounces or pounds? Rescaling a variable changes the \emph{numbers}
    OLS reports --- coefficients, standard errors, $t$, $F$, confidence intervals ---
    yet should change \textbf{nothing essential}. Often we rescale purely for cosmetics
    (fewer leading zeros).
  \end{highlightbox}

  \begin{keybox}
    \textbf{Setup.} Track how each statistic transforms when $y$ or $x_j$ is multiplied
    by a constant, and confirm that measured effects and test outcomes are preserved.
    Then rescale to \textbf{standard-deviation units} --- \emph{beta (standardized)
    coefficients} --- to compare regressors on a common footing. \textbf{Plan:}
    rescaling $y$, rescaling $x_j$, log invariance, beta coefficients.
  \end{keybox}

\end{frame}

\begin{frame}[shrink=8]{Effects of Rescaling the Dependent Variable}

  Will switching birth weight from \emph{ounces} to \emph{pounds} change any conclusion?
  Rescale $y$ and watch every statistic. Consider the birth weight equation
  \[
    \widehat{bwght} = \underset{(1.049)}{116.974}
      - \underset{(0.0916)}{0.4634}\, cigs
      + \underset{(0.0292)}{0.0927}\, faminc ,
  \]
  and define $bwghtlbs = bwght/16$ (ounces $\to$ pounds).

  \medskip

  \textbf{What changes?} Every coefficient and standard error is divided by $16$, so the
  $t$ statistics and $\Rsq$ are \textbf{unchanged}; $\SSR$ falls by $16^2 = 256$ and the
  SER by $16$.

  \begin{keybox}
    \textbf{Rule --- rescaling $y$ by $c$:}\;
    $\bhat_j \to c\,\bhat_j$ and $\se(\bhat_j) \to c\,\se(\bhat_j)$ (so $t_{\bhat_j}$ and
    $\Rsq$ are unchanged); $\SSR \to c^2\SSR$, $\;\text{SER} \to |c|\,\text{SER}$.

    \smallskip
    \textbf{Key insight:} rescaling $y$ is purely cosmetic --- it changes the appearance,
    not the conclusions.
  \end{keybox}

\end{frame}

\begin{frame}[shrink=8]{Effects of Rescaling an Independent Variable}

  Now define $packs = cigs/20$ (cigarettes $\to$ packs per day). Rewriting the equation,
  \[
    \widehat{bwght} = \underset{(1.049)}{116.974}
      - \underset{(1.832)}{9.268}\, packs
      + \underset{(0.0292)}{0.0927}\, faminc .
  \]

  \medskip

  \textbf{What changes?} Only the $packs$ terms: its coefficient becomes
  $20\times(-0.4634) = -9.268$ and its standard error $20\times0.0916 = 1.832$, leaving
  the $t$ statistic $-9.268/1.832 = -5.06$ \textbf{unchanged}. The intercept, the other
  coefficients, $\Rsq$, $\SSR$, and SER are all unchanged.

  \begin{keybox}
    \textbf{Rule --- rescaling $x_j$ by $c$:}\;
    $\bhat_j \to \bhat_j/c$ and $\se(\bhat_j) \to \se(\bhat_j)/c$ (so $t_{\bhat_j}$ is
    unchanged); the other coefficients, $\Rsq$, $\SSR$, and SER are all unchanged.
  \end{keybox}

  \medskip

  \textbf{Bottom line.} Rescaling any variable changes coefficients and standard errors
  proportionally, but never the $t$ statistics, $F$ statistics, $\Rsq$, or any test
  conclusion.

\end{frame}

\begin{frame}[shrink=8]{Rescaling Summary and Logarithmic Invariance}

      {\small
      \begin{tabular}{lcc}
        \toprule
        & \textbf{Rescale $y$} & \textbf{Rescale $x_j$} \\
        & ($y \to cy$) & ($x_j \to cx_j$) \\
        \midrule
        $\bhat_j$    & $\times c$ & $\div c$ \\
        $\se(\bhat_j)$ & $\times c$ & $\div c$ \\
        $t_{\bhat_j}$  & unchanged & unchanged \\
        $\Rsq$       & unchanged & unchanged \\
        $\SSR$       & $\times c^2$ & unchanged \\
        SER          & $\times |c|$ & unchanged \\
        \bottomrule
      \end{tabular}
      }
  \medskip

      \begin{keybox}
        \textbf{Logarithmic invariance:} \\[4pt]
        If $y$ appears as $\log(y)$: changing the units of $y$
        only affects the intercept. \\[4pt]
        $\log(c \cdot y) = \log(c) + \log(y)$ \\[4pt]
        All slope coefficients are \textbf{unchanged}. \\[6pt]
        Similarly, if $x_j$ appears as $\log(x_j)$: changing
        the units of $x_j$ only shifts the intercept. \\[4pt]
        \textbf{Elasticities are invariant to units of measurement.}
      \end{keybox}

  \medskip

  \medskip
  \textbf{Practical advice.} Use rescaling to make coefficients easy to read ---
  eliminate long strings of zeros (if $\bhat = 0.00013$, measure $x$ in thousands).
  This is cosmetic, not substantive.

\end{frame}

\begin{frame}{Beta (Standardized) Coefficients: Definition}

  When regressors live on wildly different scales --- dollars, years, test points ---
  the raw $\bhat_j$ cannot be ranked by ``importance'': a small coefficient may merely
  reflect large units. \textbf{Standardizing} puts every variable in common units.

  \begin{definitionbox}{Beta Coefficients}
    \textbf{Standardize} all variables to $z$-scores:
    \[
      z_y = \frac{y_i - \bar{y}}{\hat\sigma_y}, \qquad
      z_j = \frac{x_{ij} - \bar{x}_j}{\hat\sigma_j} .
    \]
    Then run the (intercept-free) regression
    $\hat{z}_y = \hat{b}_1 z_1 + \cdots + \hat{b}_k z_k$, whose \textbf{beta coefficients}
    are
    \[
      \hat{b}_j = \frac{\hat\sigma_j}{\hat\sigma_y}\, \hat\beta_j,
      \qquad j = 1, \ldots, k .
    \]
  \end{definitionbox}

\end{frame}

\begin{frame}{Beta Coefficients: Where the Formula Comes From}

  \begin{keybox}
    Write the fitted equation in deviations,
    $\hat y_i-\bar y=\sum_j \bhat_j(x_{ij}-\bar x_j)$; divide by $\hat\sigma_y$ and
    multiply/divide term $j$ by $\hat\sigma_j$:
    \[
      \frac{\hat y_i-\bar y}{\hat\sigma_y}
        = \sum_j \underbrace{\frac{\hat\sigma_j}{\hat\sigma_y}\,\bhat_j}_{\hat b_j}\;
          \frac{x_{ij}-\bar x_j}{\hat\sigma_j} .
    \]
    So $\hat b_j$ is simply $\bhat_j$ rescaled by the ratio of standard deviations.
  \end{keybox}

  \medskip

  \begin{highlightbox}
    \textbf{Interpretation.} A one \emph{standard-deviation} increase in $x_j$ changes
    $\hat y$ by $\hat b_j$ standard deviations --- a unit-free measure that lets you
    compare regressors directly.
  \end{highlightbox}

\end{frame}

\begin{frame}{Beta Coefficients: When Are They Useful?}

  \begin{keybox}
    \textbf{Use beta coefficients when:}
    \begin{itemize}
      \item variables are measured on very different scales (test scores, income, age);
      \item you want to rank regressors by ``importance'';
      \item the units are arbitrary or hard to interpret.
    \end{itemize}
  \end{keybox}

\end{frame}

\begin{frame}{Beta Coefficients: Key Facts}

  \begin{highlightbox}
    \begin{itemize}
      \item The $t$ statistics are \textbf{identical} for standardized and
            unstandardized variables.
      \item With a single regressor, the beta coefficient equals the sample correlation
            $r_{xy}$ (between $-1$ and $1$).
      \item Even when elasticities are the main interest, beta coefficients can still be
            informative.
    \end{itemize}
  \end{highlightbox}

\end{frame}

\begin{frame}{Example 6.1 --- Pollution and Housing Prices: The Model}

  \begin{examplebox}{Beta Coefficients for Housing Prices (HPRICE2)}
    The level--level model is
    \[
      price = \beta_0 + \beta_1 nox + \beta_2 crime + \beta_3 rooms
            + \beta_4 dist + \beta_5 stratio + u .
    \]
    Replacing each variable by its $z$-score gives the standardized estimates
    \[
      \widehat{zprice} = -0.340\, znox - 0.143\, zcrime + 0.514\, zrooms
                        - 0.235\, zdist - 0.270\, zstratio .
    \]
  \end{examplebox}

\end{frame}

\begin{frame}{Example 6.1 --- Ranking Regressors by Importance}

  \textbf{Ranking by $|\hat b|$:}
  \begin{enumerate}
    \item $rooms$: $0.514$ (largest)
    \item $nox$: $0.340$
    \item $stratio$: $0.270$
    \item $dist$: $0.235$
    \item $crime$: $0.143$ (smallest)
  \end{enumerate}

  \medskip

  \begin{keybox}
    A one-standard-deviation increase in $nox$ lowers price by $0.34$ sd --- the same
    \emph{relative} movement in pollution affects price more than crime does.
    \textbf{You cannot read this off the raw coefficients}, because their units differ.
  \end{keybox}

\end{frame}

%% ============================================================
\section{Logarithmic Functional Forms (\S6-2a)}
%% ============================================================

\begin{frame}{Logarithmic Functional Forms: Motivation \& Setup}

  \begin{highlightbox}
    \textbf{Motivation.} Logs are everywhere in applied work: they deliver
    \textbf{constant-elasticity} and \textbf{semi-elasticity} interpretations that are
    \emph{unit-free}, they narrow the range of variables, and they dampen the influence
    of outliers and heteroskedasticity --- but they make no sense for variables that can
    be zero, negative, or are already proportions.
  \end{highlightbox}

  \begin{keybox}
    \textbf{Setup.} Compare the four forms (level--level, level--log, log--level,
    log--log) and their slope interpretations, then sharpen the \emph{exact} percentage
    change $100[\exp(\hat\beta_1)-1]$ against the approximation $100\hat\beta_1$.
    \textbf{Plan:} the four forms, a housing-price example, exact vs.\ approximate
    percentages, and rules of thumb for when to use logs.
  \end{keybox}

\end{frame}

\begin{frame}[shrink=8]{Review: Four Functional Forms}

  \begin{highlightbox}
    \textbf{The four functional form combinations for $y$ and $x$:}
  \end{highlightbox}

  \medskip

  \begin{tabular}{llll}
    \toprule
    \textbf{Model} & \textbf{Dependent} & \textbf{Independent} & \textbf{Interpretation of $\beta_1$} \\
    \midrule
    Level--Level & $y$        & $x$        & $\Delta y = \beta_1 \Delta x$ \\[3pt]
    Level--Log   & $y$        & $\log(x)$  & $\Delta y = (\beta_1/100)\, \%\Delta x$ \\[3pt]
    Log--Level   & $\log(y)$  & $x$        & $\%\Delta y \approx 100\beta_1 \Delta x$ \\[3pt]
    Log--Log     & $\log(y)$  & $\log(x)$  & $\%\Delta y \approx \beta_1\, \%\Delta x$ (elasticity) \\
    \bottomrule
  \end{tabular}

  \medskip

  \begin{keybox}
    \textbf{Remember:} The log-level and log-log interpretations are
    \emph{approximations}. For large changes, use the exact formula
    $\%\Delta y = 100[\exp(\bhat \Delta x) - 1]$ (next slide).
  \end{keybox}

\end{frame}

\begin{frame}[shrink=8]{Functional Forms in Practice: Housing Prices}

  \begin{examplebox}{Housing Prices (HPRICE2 --- eq.\ 6.7)}
    \[
      \widehat{\log(price)} = \underset{(0.19)}{9.23}
        - \underset{(0.066)}{0.718}\, \log(nox)
        + \underset{(0.019)}{0.306}\, rooms
    \]
    \[
      n = 506, \quad \Rsq = 0.514
    \]
  \end{examplebox}

  \medskip

      \textbf{Interpreting the coefficients:}
      \begin{itemize}
        \item $\bhat_{\log(nox)} = -0.718$: elasticity \\
              1\% increase in $nox$ $\Rightarrow$ price falls by $\approx 0.718\%$
        \item $\bhat_{rooms} = 0.306$: semi-elasticity \\
              One more room $\Rightarrow$ price rises by $\approx 30.6\%$
      \end{itemize}
  \medskip

      \begin{keybox}
        This equation mixes log--log (for $nox$) and log--level (for $rooms$). \\[4pt]
        Both interpretations follow directly from the functional form table. \\[4pt]
        The 30.6\% approximation for $rooms$ is somewhat rough ---
        the exact change is 35.8\% (next slide).
      \end{keybox}

\end{frame}

\begin{frame}{Exact vs.\ Approximate Percentage Changes}

  \begin{definitionbox}{Exact Percentage Change (eq.\ 6.8)}
    In the model $\widehat{\log(y)} = \bhat_0 + \bhat_1 x_1 + \bhat_2 x_2$, fixing $x_1$,
    \[
      \%\Delta\hat{y} = 100 \cdot [\exp(\bhat_2 \Delta x_2) - 1] .
    \]
  \end{definitionbox}

  \medskip

  \textbf{When does this matter?} For small $\bhat_2\Delta x_2$ the approximation
  $100\,\bhat_2\Delta x_2$ is fine; for large changes the exact formula differs
  substantially.

  \medskip

  \textbf{Why?} A first-order Taylor expansion $\exp(a)\approx 1+a$ gives
  $100\,\bhat_2\Delta x_2$ --- only the first-order term; the dropped $a^2/2$ is why the
  exact increase and decrease are asymmetric.

\end{frame}

\begin{frame}{Exact vs.\ Approximate: Housing Example}

  \begin{examplebox}{Housing Prices: the $rooms$ coefficient}
    Here $\bhat_{rooms} = 0.306$.

    \smallskip
    \textbf{Add one room} ($\Delta rooms = 1$): approx $30.6\%$, exact
    $100[\exp(0.306) - 1] = 35.8\%$.

    \smallskip
    \textbf{Remove one room} ($\Delta rooms = -1$): approx $-30.6\%$, exact
    $100[\exp(-0.306) - 1] = -26.4\%$.
  \end{examplebox}

  \medskip

  \begin{keybox}
    The approximation always lies \emph{between} the exact increase and the exact
    decrease --- and the gap widens as $|\bhat_2\Delta x_2|$ grows.
  \end{keybox}

\end{frame}

\begin{frame}{When to Use Logarithms vs.\ Levels}

  \textbf{Use $\log(x)$ when:}
  \begin{itemize}
    \item $x > 0$ always (wages, prices, sales, population);
    \item $x$ varies a lot (firm sales span millions to billions);
    \item you want a percentage-change interpretation, less outlier influence, or more
      plausible normality.
  \end{itemize}

  \medskip

  \textbf{Use $x$ in levels when:}
  \begin{itemize}
    \item $x$ is in years (education, experience, age);
    \item $x$ is a proportion or percent (unemployment rate, pass rate);
    \item $x$ can be zero or negative.
  \end{itemize}

\end{frame}

\begin{frame}{Logarithms: Benefits and the Zero-Value Caveat}

  \begin{keybox}
    \textbf{Benefits of logs:}
    \begin{enumerate}
      \item coefficients carry a percentage-change interpretation;
      \item slope coefficients are invariant to units;
      \item the range of the variable narrows $\Rightarrow$ less outlier sensitivity;
      \item $\log(y)$ often makes the CLM assumptions (homoskedasticity, normality) more
        plausible.
    \end{enumerate}
  \end{keybox}

  \medskip

  \begin{highlightbox}
    \textbf{Caveat.} If $y \ge 0$ with $y = 0$ possible, $\log(y)$ is undefined. The fix
    $\log(1+y)$ is sometimes used, but its coefficients are \textbf{not} invariant to the
    units of $y$ --- use with caution.
  \end{highlightbox}

\end{frame}

%% ============================================================
\section{Quadratic Models and Interactions (\S6-2b, \S6-2c)}
%% ============================================================

\begin{frame}{Quadratics and Interactions: Motivation \& Setup}

  \begin{highlightbox}
    \textbf{Motivation.} A linear model forces a \emph{constant} partial effect, but
    economics often demands otherwise: the return to experience \textbf{diminishes},
    and the payoff to class attendance may \textbf{depend on} ability. Quadratics and
    interactions relax the constant-effect straitjacket while staying linear in the
    parameters.
  \end{highlightbox}

  \begin{keybox}
    \textbf{Setup.} The partial effect is now itself a \emph{function}, found by
    calculus: a quadratic gives $\partial y/\partial x = \beta_1 + 2\beta_2 x$, and an
    interaction gives $\partial y/\partial x_1 = \beta_1 + \beta_3 x_2$. \textbf{Plan:}
    quadratics and turning points, interaction terms, and worked examples.
  \end{keybox}

\end{frame}

\begin{frame}{Models with Quadratics}

  The simplest way to let a marginal effect \emph{change} with the level of $x$ ---
  diminishing returns to experience, a U-shaped cost curve --- is to add the square
  of $x$ as its own regressor.

  \begin{definitionbox}{Quadratic Specification}
    \[
      y = \beta_0 + \beta_1 x + \beta_2 x^2 + u, \qquad
      \hat{y} = \bhat_0 + \bhat_1 x + \bhat_2 x^2 .
    \]
    The \textbf{marginal effect} of $x$ now depends on $x$:
    \[
      \frac{\Delta \hat{y}}{\Delta x} \approx \bhat_1 + 2\bhat_2 x .
    \]
  \end{definitionbox}

\end{frame}

\begin{frame}{Quadratics: Turning Point and Shape}

  \begin{keybox}
    \textbf{Turning point.} Set the marginal effect to zero,
    $\bhat_1 + 2\bhat_2 x^* = 0$:
    \[
      x^* = -\frac{\bhat_1}{2\bhat_2} = \frac{|\bhat_1|}{2|\bhat_2|} .
    \]
    Always check whether $x^*$ falls within the range of the data.
  \end{keybox}

  \medskip

  \textbf{Two shapes:}
  \begin{itemize}
    \item $\bhat_1 > 0,\ \bhat_2 < 0$: \textbf{hump} (diminishing returns);
    \item $\bhat_1 < 0,\ \bhat_2 > 0$: \textbf{U-shape} (increasing returns);
    \item same signs on $\bhat_1,\bhat_2$: monotonic, no turning point for $x > 0$.
  \end{itemize}

\end{frame}

\begin{frame}{Wage--Experience Quadratic (WAGE1)}

  \begin{examplebox}{Diminishing Returns to Experience (eq.\ 6.12)}
    \[
      \widehat{wage} = \underset{(0.35)}{3.73}
        + \underset{(0.041)}{0.298}\, exper
        - \underset{(0.0009)}{0.0061}\, exper^2
    \]
    \[
      n = 526, \quad \Rsq = 0.093 .
    \]
  \end{examplebox}

  \medskip

  The marginal effect
  $\dfrac{\Delta \widehat{wage}}{\Delta exper} \approx 0.298 - 2(0.0061)\, exper$
  \emph{falls} as experience rises.

\end{frame}

\begin{frame}{Wage--Experience: Diminishing Returns}

  \textbf{Marginal effect at different experience levels:}
  \begin{itemize}
    \item $exper = 1$: $\ 0.298 - 0.012 = 0.286$ (\$0.29/hr)
    \item $exper = 10$: $\ 0.298 - 0.122 = 0.176$ (\$0.18/hr)
    \item $exper = 20$: $\ 0.298 - 0.244 = 0.054$ (\$0.05/hr)
  \end{itemize}

  \medskip

  \begin{keybox}
    \textbf{Turning point:}
    $exper^* = \dfrac{0.298}{2(0.0061)} \approx 24.4$ years. After that the estimated
    return to experience turns \textbf{negative}. About $28\%$ of the sample has
    $exper > 24$, so the turning point is economically relevant --- not a mathematical
    artifact.
  \end{keybox}

\end{frame}

\begin{frame}[shrink=8]{Example 6.2 --- Housing Prices with Quadratic $rooms$}

  \begin{examplebox}{Quadratic in $rooms$ (HPRICE2 --- eq.\ 6.14)}
    \[
      \widehat{\log(price)} = \underset{(0.57)}{13.39}
        - \underset{(0.115)}{0.902}\, \log(nox)
        - \underset{(0.043)}{0.087}\, \log(dist)
    \]
    \[
      - \underset{(0.165)}{0.545}\, rooms
      + \underset{(0.013)}{0.062}\, rooms^2
      - \underset{(0.006)}{0.048}\, stratio
    \]
    \[
      n = 506, \quad \Rsq = 0.603
    \]
  \end{examplebox}

  \medskip

  \textbf{The U-shape:} $\bhat_1 < 0$, $\bhat_2 > 0$. \quad
  Turning point: $rooms^* = 0.545 / [2(0.062)] \approx 4.4$.
  Only 1\% of communities have $\leq 4.4$ average rooms,
  so for the vast majority more rooms $\Rightarrow$ higher price.

\end{frame}

\begin{frame}[shrink=8]{Example 6.2 --- Marginal Effect of $rooms$}

  \begin{keybox}
    \textbf{Marginal effect of $rooms$ on $\log(price)$:}
    \[
      \%\Delta \widehat{price} \approx (-54.5 + 12.4\, rooms)\, \Delta rooms
    \]
  \end{keybox}

  \medskip

      \textbf{Effect at different values:}
      \begin{itemize}
        \item At $rooms = 5$: $(-54.5 + 62.0) = 7.5\%$ per room
        \item At $rooms = 6$: $(-54.5 + 74.4) = 19.9\%$ per room
        \item At $rooms = 7$: $(-54.5 + 86.8) = 32.3\%$ per room
      \end{itemize}

      \medskip

      A very strong \textbf{increasing} effect at typical values!
  \medskip

      \begin{highlightbox}
        \textbf{Lesson:} The coefficient on $rooms^2$ is just 0.062 ---
        seemingly tiny. But the marginal effect $\bhat_1 + 2\bhat_2 \cdot rooms$
        changes rapidly across the data range. \\[4pt]
        Never dismiss a quadratic term based on the size
        of the coefficient alone. Always compute the marginal
        effect at relevant values of $x$.
      \end{highlightbox}

\end{frame}

\begin{frame}{Important Lesson: Small Coefficients on $x^2$ Can Matter}

  \begin{highlightbox}
    \textbf{Do not dismiss a quadratic term just because $\bhat_2$ ``looks small.''}
  \end{highlightbox}

  \medskip

  In the housing example $\bhat_{rooms^2} = 0.062$ --- seemingly tiny --- yet the
  marginal effect $\bhat_1 + 2\bhat_2\cdot rooms$ is $0.255$ ($25.5\%$/room) at
  $rooms = 6.45$ and $0.323$ ($32.3\%$/room) at $rooms = 7$. The term is also highly
  significant: $t = 0.062/0.013 = 4.77$.

\end{frame}

\begin{frame}{Why the Quadratic Beats the Linear Model Here}

  \begin{keybox}
    \textbf{Compare with a linear model.} Dropping $rooms^2$, the coefficient on $rooms$
    becomes about $0.255$ --- a flat $25.5\%$/room regardless of house size. The
    quadratic reveals the effect is \textbf{much stronger for larger houses}, an insight
    the linear model misses entirely.
  \end{keybox}

  \medskip

  \textbf{General rule:} always compute the marginal effect at interesting values of $x$
  (mean, quartiles) --- never judge a quadratic by the size of $\bhat_2$ alone.

\end{frame}

\begin{frame}[shrink=8]{Interaction Terms: When One Effect Depends on Another}

  Sometimes the effect of $x_2$ is not constant but \textbf{depends on the level of
  $x_1$} --- an extra bedroom is worth more in a larger house, and class attendance
  helps stronger students more. An \emph{interaction term} $x_1 x_2$ captures this.

  \begin{definitionbox}{Interaction model and its partial effect}
    \[
      y = \beta_0 + \beta_1 x_1 + \beta_2 x_2 + \beta_3\, x_1 x_2 + u
      \quad\Longrightarrow\quad
      \frac{\partial y}{\partial x_2} = \beta_2 + \beta_3 x_1 .
    \]
    The effect of $x_2$ now moves with $x_1$ (and symmetrically,
    $\partial y/\partial x_1 = \beta_1 + \beta_3 x_2$).
  \end{definitionbox}

  \begin{keybox}
    \textbf{Reading $\beta_2$ correctly.} $\beta_2$ is the effect of $x_2$ \emph{only
    when $x_1 = 0$} --- often uninteresting or impossible (a house with zero square
    feet). Report the partial effect at \textbf{meaningful} values of $x_1$ --- the
    mean, median, or quartiles --- instead.
  \end{keybox}

\end{frame}

\begin{frame}[shrink=8]{Example 6.3 --- Attendance and Exam Performance: Model}

  \begin{examplebox}{Interactions with Quadratics (ATTEND --- eq.\ 6.18--6.19)}
    \[
      \begin{aligned}
        stndfnl = {}& \beta_0 + \beta_1\, atndrte + \beta_2\, priGPA + \beta_3\, ACT \\
          &+ \beta_4\, priGPA^2 + \beta_5\, ACT^2 + \beta_6\, priGPA \cdot atndrte + u
      \end{aligned}
    \]
    \[
      \widehat{stndfnl} = \underset{(1.36)}{2.05}
        - \underset{(0.0102)}{0.0067}\, atndrte
        - \underset{(0.48)}{1.63}\, priGPA
        - \underset{(0.098)}{0.128}\, ACT
    \]
    \[
      + \underset{(0.101)}{0.296}\, priGPA^2
      + \underset{(0.0022)}{0.0045}\, ACT^2
      + \underset{(0.0043)}{0.0056}\, priGPA \cdot atndrte
    \]
    \[
      n = 680, \quad \Rsq = 0.229, \quad \aRsq = 0.222
    \]
  \end{examplebox}

\end{frame}

\begin{frame}[shrink=8]{Example 6.3 --- Partial Effect of Attendance}

  \begin{keybox}
    \textbf{Partial effect of attendance on standardized final exam score:}
    \[
      \frac{\Delta stndfnl}{\Delta atndrte} = \bhat_1 + \bhat_6\, priGPA
      = -0.0067 + 0.0056 \cdot priGPA
    \]
  \end{keybox}

  \medskip

      \textbf{At different prior GPA levels:}
      \begin{itemize}
        \item $priGPA = 2.0$: $-0.0067 + 0.0112 = 0.0045$
        \item $priGPA = 2.59$ (mean): $-0.0067 + 0.0145 = 0.0078$
        \item $priGPA = 3.5$: $-0.0067 + 0.0196 = 0.0129$
      \end{itemize}

      \medskip

      At the mean: a 10-point increase in attendance
      raises the standardized score by 0.078 sd.
  \medskip

      \begin{highlightbox}
        \textbf{Key insight:} \\
        The effect of attendance is \textbf{stronger} for students
        with higher prior GPA. \\[4pt]
        Students who were already doing well benefit more
        from going to class --- the interaction $priGPA \cdot atndrte$
        captures this heterogeneity.
      \end{highlightbox}

\end{frame}

\begin{frame}[shrink=8]{Interpreting the Attendance Interaction}

      \textbf{Testing the partial effect at the mean:}

      \medskip

      To get $\se$ of the effect at $\overline{priGPA} = 2.59$:
      \begin{enumerate}
        \item Replace $priGPA \cdot atndrte$ with $(priGPA - 2.59) \cdot atndrte$
        \item Re-run the regression
        \item The new coefficient on $atndrte$ equals the partial effect at the mean,
              with its own $\se$
      \end{enumerate}

      \medskip

      Result: $\bhat = 0.0078$, $\se = 0.0026$, $t = 3.0$. \\
      \textbf{Statistically significant} at the 1\% level.
  \medskip

      \begin{highlightbox}
        \textbf{Key lessons from this example:} \\[4pt]
        \begin{enumerate}
          \item Looking at $\bhat_1 = -0.0067$ alone would
                suggest attendance \emph{hurts} performance ---
                but this is the effect when $priGPA = 0$ (meaningless!)
          \item The individual $t$ stats on $\bhat_1$ and $\bhat_6$ are both
                insignificant, but the $F$-test for joint significance
                rejects $\text{H}_0$: $\beta_1 = 0, \beta_6 = 0$ at 5\%
          \item \textbf{Always evaluate interactions at interesting values}
        \end{enumerate}
      \end{highlightbox}

\end{frame}

%% ============================================================
\section{Average Partial Effects and Centering (\S6-2d)}
%% ============================================================

\begin{frame}{Average Partial Effects and Centering: Motivation \& Setup}

  \begin{highlightbox}
    \textbf{Motivation.} Once the partial effect varies across people (quadratics,
    interactions), \emph{which} number do we report? And why do the coefficients on the
    level terms look so strange after adding an interaction? Both puzzles are solved by
    averaging --- and by centering variables before interacting them.
  \end{highlightbox}

  \begin{keybox}
    \textbf{Setup.} The \textbf{average partial effect (APE)} averages the
    observation-specific effect over the sample,
    $\text{APE}_j = n^{-1}\sum_i \partial \hat y_i / \partial x_j$. Centering each
    variable at its mean before forming the interaction makes the coefficient on the
    level term read directly as the APE. \textbf{Plan:} defining the APE, then
    centering.
  \end{keybox}

\end{frame}

\begin{frame}{Average Partial Effects (APE)}

  \begin{definitionbox}{Average Partial Effect}
    In a model with quadratics or interactions, the partial effect of $x_j$ varies
    across observations. The \textbf{average partial effect (APE)} averages it over the
    sample:
    \[
      \widehat{\text{APE}}_j = \frac{1}{n} \sum_{i=1}^{n}
      \left. \frac{\partial \hat{y}_i}{\partial x_j} \right|_{x_i} .
    \]
  \end{definitionbox}

\end{frame}

\begin{frame}{APE: The Attendance Model}

  \begin{examplebox}{APE of attendance}
    Partial effect of $atndrte$ is $\bhat_1 + \bhat_6\, priGPA_i$, so
    \[
      \widehat{\text{APE}}_{atndrte} = \bhat_1 + \bhat_6\, \overline{priGPA}
        = -0.0067 + 0.0056(2.59) = 0.0078 .
    \]
  \end{examplebox}

  \medskip

  \begin{examplebox}{APE of prior GPA}
    Partial effect is $\bhat_2 + 2\bhat_4\, priGPA_i + \bhat_6\, atndrte_i$, so
    \[
      \widehat{\text{APE}}_{priGPA} = \bhat_2 + 2\bhat_4\, \overline{priGPA}
        + \bhat_6\, \overline{atndrte} .
    \]
  \end{examplebox}

\end{frame}

\begin{frame}{Centering Variables Before Creating Interactions}

  In $y = \beta_0 + \beta_1 x_1 + \beta_2 x_2 + \beta_3 x_1 x_2 + u$, the coefficient
  $\beta_2$ is the effect of $x_2$ \emph{only when $x_1 = 0$} --- often uninteresting.

  \medskip

  \textbf{Reparameterize:} replace $x_1 x_2$ with $(x_1 - \bar{x}_1)(x_2 - \bar{x}_2)$:
  \[
    y = \alpha_0 + \delta_1 x_1 + \delta_2 x_2
      + \beta_3 (x_1 - \bar{x}_1)(x_2 - \bar{x}_2) + u .
  \]
  Now $\delta_2 = \beta_2 + \beta_3 \bar{x}_1$ is exactly the \textbf{APE} of $x_2$.

\end{frame}

\begin{frame}{Centering: Benefits}

  \begin{keybox}
    \textbf{Benefits of centering:}
    \begin{enumerate}
      \item the coefficients on $x_1$ and $x_2$ become the APE directly;
      \item standard errors for the APE come automatically;
      \item it reduces multicollinearity between $x_j$ and $x_j x_k$;
      \item the interaction coefficient $\beta_3$ is unchanged.
    \end{enumerate}
  \end{keybox}

  \medskip

  \begin{highlightbox}
    The same logic applies to quadratics: in
    $y = \gamma_0 + \gamma_1 x + \beta_2(x - \bar{x})^2 + u$, the coefficient $\gamma_1$
    is the APE of $x$.
  \end{highlightbox}

\end{frame}

%% ============================================================
\section{Adjusted $\Rsq$ and Model Selection (\S6-3)}
%% ============================================================

\begin{frame}{Adjusted $\Rsq$ and Model Selection: Motivation \& Setup}

  \begin{highlightbox}
    \textbf{Motivation.} $R^2$ \emph{never} falls when you add a regressor --- so it
    cannot tell you whether a variable belongs. We need a measure that
    \textbf{penalizes} complexity, a principled way to choose between competing models,
    and the judgment to know when \emph{not} to add a control.
  \end{highlightbox}

  \begin{keybox}
    \textbf{Setup.} The \textbf{adjusted} $\Rsq$ replaces sample sums with
    degrees-of-freedom-corrected estimators,
    $\aRsq = 1 - \dfrac{\hat\sigma^2}{\text{SST}/(n-1)}$, so it can \emph{fall} when a
    useless regressor is added. \textbf{Plan:} $\aRsq$ and its properties, choosing
    between nonnested models, over-controlling, and adding regressors to cut the error
    variance.
  \end{keybox}

\end{frame}

\begin{frame}[shrink=8]{Adjusted $R$-Squared: Definition}

  \textbf{Recall:} $\Rsq = 1 - \SSR/\SST$ never falls when a variable is added. \\
  To penalize model complexity, replace $\SSR/n$ with $\SSR/(n-k-1)$:

  \medskip

  \begin{definitionbox}{Adjusted $R$-Squared (eq.\ 6.21)}
    \[
      \aRsq = 1 - \frac{\SSR/(n-k-1)}{\SST/(n-1)}
            = 1 - \frac{\hat\sigma^2}{\SST/(n-1)}
    \]
    Alternative formula: $\displaystyle\aRsq = 1 - (1 - \Rsq)\frac{n-1}{n-k-1}$
  \end{definitionbox}

  \medskip

  \begin{keybox}
    \textbf{Why $n-k-1$?} \\[4pt]
    $\Rsq = 1 - \SSR/n \div \SST/n$. Replacing with unbiased estimators:
    $\hat\sigma^2 = \SSR/(n-k-1)$ and $\SST/(n-1)$ is the sample variance of $y$.
    Using unbiased estimators instead of $n$-scaled versions gives $\aRsq$.
  \end{keybox}

\end{frame}

\begin{frame}[shrink=8]{Adjusted $R$-Squared: Key Properties}

  \begin{keybox}
    \textbf{Properties of $\aRsq$:}
    \begin{enumerate}
      \item Penalizes adding variables: $\aRsq$ \textbf{can decrease} when a variable is added
      \item $\aRsq$ increases $\iff$ the $t$ statistic on the new variable exceeds 1 in absolute value
      \item For a group of variables: $\aRsq$ increases $\iff$ the $F$ statistic $> 1$
      \item $\aRsq \leq \Rsq$ always; with small $n$ and large $k$, $\aRsq$ can be \textbf{negative}
      \item $\aRsq < 0$ signals a very poor model relative to its complexity
    \end{enumerate}
  \end{keybox}

  \medskip

  \begin{highlightbox}
    \textbf{Important:} The $F$ statistic for testing exclusion restrictions uses $\Rsq$, \textbf{not} $\aRsq$.
    The formula with $\aRsq$ is \emph{not valid}!
  \end{highlightbox}

\end{frame}

\begin{frame}[shrink=8]{$\aRsq$ Can Be Negative}

      \begin{examplebox}{Numerical Example}
        If $\Rsq = 0.10$, $n = 51$, $k = 10$:
        \[
          \aRsq = 1 - (1 - 0.10) \frac{50}{40} = 1 - 1.125 = -0.125
        \]
        A negative $\aRsq$ tells us the model has too many
        regressors relative to the sample size and explanatory power.
      \end{examplebox}
  \medskip

      \begin{highlightbox}
        \textbf{Important reminders:}
        \begin{itemize}
          \item The $F$ statistic in (4.41) uses $\Rsq$, \textbf{not} $\aRsq$.
                Using $\aRsq$ in the $F$ formula is \emph{not} valid!
          \item When $\aRsq$ is reported alongside $\Rsq$, remember that sometimes
                the $\aRsq$ column is mislabeled as $\Rsq$ --- always check
          \item $\aRsq$ is most useful for comparing models with the
                \textbf{same dependent variable} but different regressors
        \end{itemize}
      \end{highlightbox}

\end{frame}

\begin{frame}{Choosing Between Nonnested Models}

  \textbf{Nested models} (one is a special case of the other) $\Rightarrow$ use the
  $F$ test. \textbf{Nonnested models} (neither is a special case) $\Rightarrow$ the
  $F$ test does not apply.

  \medskip

  \begin{examplebox}{Baseball Salary (MLB1)}
    Model A: $\log(salary)$ on $years$, $gamesyr$, $bavg$, $hrunsyr$. \\
    Model B: same, but $rbisyr$ instead of $hrunsyr$.

    \smallskip
    $\aRsq_A = 0.6211$, \quad $\aRsq_B = 0.6226$ --- a slight, tiny preference for B.
  \end{examplebox}

\end{frame}

\begin{frame}{Comparing Nonnested Models with $\aRsq$}

  \begin{keybox}
    \textbf{Pick the model with the higher $\aRsq$} --- valid only when the models share
    the \textbf{same dependent variable} (same $\SST$) and the \textbf{same sample}
    (same $n$).
  \end{keybox}

  \medskip

  \begin{highlightbox}
    \textbf{Critical limitation.} You cannot use $\Rsq$ or $\aRsq$ to compare $y$ vs.\
    $\log(y)$ as dependent variables --- they measure variation in \emph{different}
    things. Section~6-4 gives the fix.
  \end{highlightbox}

\end{frame}

\begin{frame}[shrink=8]{Example 6.4 --- CEO Compensation: Level vs.\ Log}

  \begin{examplebox}{Two Models for CEO Salary (CEOSAL2)}
    \textbf{Model A} (level--level):
    \[
      \widehat{salary} = \underset{(223.90)}{830.63}
        + \underset{(0.0089)}{0.0163}\, sales
        + \underset{(11.08)}{19.63}\, roe
    \]
    \[
      n = 209, \quad \Rsq = 0.029, \quad \aRsq = 0.020
    \]

    \medskip

    \textbf{Model B} (log--log):
    \[
      \widehat{lsalary} = \underset{(0.29)}{4.36}
        + \underset{(0.033)}{0.275}\, lsales
        + \underset{(0.0040)}{0.0179}\, roe
    \]
    \[
      n = 209, \quad \Rsq = 0.282, \quad \aRsq = 0.275
    \]
  \end{examplebox}

  \medskip

\end{frame}

\begin{frame}{Example 6.4 --- Why You Cannot Compare the $\Rsq$}

  \begin{keybox}
    You cannot compare $\Rsq = 0.029$ (for $salary$) with $\Rsq = 0.282$ (for
    $\log(salary)$) --- they explain variation in \textbf{different} dependent variables.
    A valid comparison needs the same-scale pseudo-$\Rsq$ from Section~6-4.
  \end{keybox}

\end{frame}

\begin{frame}{Over-Controlling: Adding Bad Controls}

  \textbf{It is possible to control for \emph{too many} variables:} including a
  ``bad control'' can introduce bias or mask the very effect you want to measure.

  \medskip

  \begin{examplebox}{Beer Taxes and Traffic Fatalities}
    Goal: the effect of beer $tax$ on $fatalities$.

    \smallskip
    \textbf{Good controls:} $miles$, $percmale$, $perc16\_21$ (demographics, exposure).

    \smallskip
    \textbf{Bad control:} $beercons$. Holding beer consumption fixed, $tax$ could affect
    fatalities only through non-alcohol channels --- not what we want.
  \end{examplebox}

\end{frame}

\begin{frame}{When Is a Control ``Bad''?}

  \begin{keybox}
    A variable is a \textbf{bad control} if it is
    \begin{enumerate}
      \item on the \textbf{causal pathway} between $x$ and $y$ (a ``mediator'');
      \item an \textbf{outcome} of the treatment variable;
      \item one whose inclusion changes the ceteris paribus interpretation into
        something uninteresting.
    \end{enumerate}

    \smallskip
    \textbf{Rule of thumb:} ask whether holding this variable fixed makes economic sense.
  \end{keybox}

\end{frame}

\begin{frame}{Adding Regressors to Reduce Error Variance}

  \textbf{The trade-off.} Adding a variable to a regression:
  \begin{itemize}
    \item[\textcolor{positive}{$+$}] reduces the error variance $\sigma^2$ (takes
      variation out of $u$);
    \item[\textcolor{negative}{$-$}] may increase multicollinearity (inflates
      $\Var(\bhat_j)$).
  \end{itemize}

  \medskip

  \textbf{Clear case.} If the new variables are \textbf{uncorrelated} with the regressors
  of interest, there is no multicollinearity cost, $\sigma^2$ falls, and all standard
  errors shrink --- so \textbf{always include} them.

\end{frame}

\begin{frame}{Reducing Error Variance: A Random-Assignment Example}

  \begin{examplebox}{Computer Equipment Grants}
    Effect of a \emph{randomly assigned} grant amount on college GPA:
    \begin{itemize}
      \item random assignment $\Rightarrow$ the grant is uncorrelated with HS GPA,
        family background, etc.;
      \item adding HS GPA and test scores as controls reduces $\hat\sigma^2$ without
        increasing multicollinearity;
      \item result: a more precise estimate of the grant effect.
    \end{itemize}
  \end{examplebox}

  \medskip

  \begin{keybox}
    \textbf{With random assignment,} include pre-treatment controls to improve precision
    --- even though they are not needed for unbiasedness.
  \end{keybox}

\end{frame}

%% ============================================================
\section{Prediction and Residual Analysis (\S6-4)}
%% ============================================================

\begin{frame}{Prediction and Residual Analysis: Motivation \& Setup}

  \begin{highlightbox}
    \textbf{Motivation.} So far we used regression to estimate ceteris paribus
    \emph{effects}. But we also want to \textbf{predict} $y$ at chosen values of the
    regressors --- and to be honest about the uncertainty, which differs for an
    \emph{average} outcome versus a \emph{single} new unit.
  \end{highlightbox}

  \begin{keybox}
    \textbf{Setup.} Distinguish the \textbf{confidence interval} for the mean
    $\E(y\mid\mathbf{x}_0)$ from the wider \textbf{prediction interval} for an individual
    $y_0$ (which adds the error variance $\hat\sigma^2$). \textbf{Plan:} confidence
    intervals for predictions, prediction intervals, residual analysis, and recovering
    $y$ when $\log(y)$ is the dependent variable.
  \end{keybox}

\end{frame}

\begin{frame}{Confidence Intervals for Predictions}

  We often want the predicted \emph{mean} response at chosen values $x_j = c_j$, with a
  confidence interval. Given $\hat{y} = \bhat_0 + \bhat_1 x_1 + \cdots + \bhat_k x_k$, the
  point prediction is
  \[
    \hat\theta_0 = \bhat_0 + \bhat_1 c_1 + \cdots + \bhat_k c_k .
  \]

  \medskip

  \textbf{The reparameterization trick} (to get $\se(\hat\theta_0)$):
  \begin{enumerate}
    \item define $x_j^0 = x_j - c_j$ for each variable;
    \item regress $y$ on $x_1^0, \ldots, x_k^0$;
    \item the new intercept equals $\hat\theta_0$ --- with its own standard error.
  \end{enumerate}
  The 95\% CI is then $\hat\theta_0 \pm t_{0.025}\,\se(\hat\theta_0)$.

\end{frame}

\begin{frame}{Why the Reparameterization Works}

  \begin{keybox}
    Write $\beta_0 = \theta_0 - \beta_1 c_1 - \cdots - \beta_k c_k$ and substitute:
    \[
      y = \theta_0 + \beta_1(x_1 - c_1) + \cdots + \beta_k(x_k - c_k) + u .
    \]
    The intercept of this regression is exactly $\theta_0 = \E[y \mid x = c]$, so OLS
    reports $\hat\theta_0$ and its standard error directly.
  \end{keybox}

  \medskip

  \begin{highlightbox}
    The prediction variance is \textbf{smallest at the sample means} ($c_j = \bar{x}_j$);
    predictions far from the means have wider intervals.
  \end{highlightbox}

\end{frame}

\begin{frame}{Example 6.5 --- College GPA: The Model}

  \begin{examplebox}{College GPA Prediction (GPA2 --- eq.\ 6.32)}
    \[
      \widehat{colgpa} = \underset{(0.075)}{1.493}
        + \underset{(0.00007)}{0.00149}\, sat
        - \underset{(0.00056)}{0.01386}\, hsperc
    \]
    \[
        - \underset{(0.01650)}{0.06088}\, hsize
        + \underset{(0.00227)}{0.00546}\, hsize^2
    \]
    \[
      n = 4{,}137, \quad \Rsq = 0.278, \quad \hat\sigma = 0.560
    \]
  \end{examplebox}

\end{frame}

\begin{frame}{Example 6.5 --- Confidence Interval for the Mean}

  \textbf{Predict for} $sat = 1{,}200$, $hsperc = 30$, $hsize = 5$:
  $\widehat{colgpa} \approx 2.70$, and the reparameterization trick gives
  $\se(\hat\theta_0) = 0.020$.

  \medskip

  \begin{keybox}
    \textbf{95\% CI for expected GPA:}
    \[
      2.70 \pm 1.96(0.020) = [2.66, 2.74] .
    \]
    Very narrow --- with $n = 4{,}137$ the mean is precisely estimated. This is a CI for
    the \textbf{average} GPA of all students with these characteristics.
  \end{keybox}

\end{frame}

\begin{frame}{Prediction Intervals vs.\ Confidence Intervals}

  \begin{definitionbox}{Two Types of Intervals}
    \textbf{Confidence interval:} for $\E[y \mid \mathbf{x}]$ (the average outcome for a
    subpopulation). \\[4pt]
    \textbf{Prediction interval:} for $y^0$ (a specific individual's future outcome).
  \end{definitionbox}

  \medskip

  The prediction error is $\hat{e}^0 = y^0 - \hat{y}^0$. The new unit's error $u^0$ is
  independent of the estimation sample, so $\Cov(\hat{y}^0, u^0)=0$ and the two variances
  add:
  \[
    \Var(\hat{e}^0) = \Var(\hat{y}^0) + \sigma^2 .
  \]
  \textbf{Two sources of uncertainty:} sampling error in $\bhat_j$ ($\to 0$ as
  $n\to\infty$), and the error variance $\sigma^2$ (which does \textbf{not} shrink with
  $n$).

\end{frame}

\begin{frame}{Prediction Intervals: Standard Error and Width}

  \begin{keybox}
    \textbf{Standard error of prediction:}
    $\se(\hat{e}^0) = \sqrt{[\se(\hat{y}^0)]^2 + \hat\sigma^2}$.

    \smallskip
    \textbf{95\% prediction interval:}
    $\hat{y}^0 \pm t_{0.025}\,\se(\hat{e}^0)$.
  \end{keybox}

  \medskip

  \begin{highlightbox}
    The prediction interval is \textbf{always wider} than the confidence interval because
    it includes $\hat\sigma^2$. For large $n$, $\se(\hat{y}^0)$ is tiny, so
    $\se(\hat{e}^0) \approx \hat\sigma$ --- the width is dominated by the irreducible
    error variance.
  \end{highlightbox}

\end{frame}

\begin{frame}[shrink=8]{Example 6.6 --- Prediction Interval for Future College GPA}

  \begin{examplebox}{Individual Prediction (GPA2)}
    Same student: $sat = 1{,}200$, $hsperc = 30$, $hsize = 5$. \\[4pt]
    $\hat{y}^0 = 2.70$, $\se(\hat{y}^0) = 0.020$, $\hat\sigma = 0.560$
    \[
      \se(\hat{e}^0) = \sqrt{(0.020)^2 + (0.560)^2} \approx 0.560
    \]
    \textbf{95\% prediction interval:} $2.70 \pm 1.96(0.560) = [1.60, 3.80]$
  \end{examplebox}

  \medskip

      \textbf{Compare:}
      \begin{itemize}
        \item CI for average GPA: $[2.66, 2.74]$ (very narrow)
        \item PI for individual GPA: $[1.60, 3.80]$ (very wide!)
      \end{itemize}

      \medskip

      The prediction interval spans more than 2 full GPA points.

\end{frame}

\begin{frame}{Example 6.6 --- Why the Prediction Interval Is So Wide}

  \begin{keybox}
    The model explains only $27.8\%$ of the variation in college GPA ($\Rsq = 0.278$);
    unobserved factors (motivation, study habits, health, environment) account for most
    of it.

    \smallskip
    \textbf{Good news:} SAT and HS rank do not predetermine college success.
  \end{keybox}

\end{frame}

\begin{frame}{Residual Analysis}

  The residual itself carries information: a large $|\uh_i|$ flags an observation the
  model badly misses --- an underpriced house, a value-adding school, or an over- or
  under-paid athlete.

  \begin{definitionbox}{Residual Analysis}
    Examine individual residuals $\uh_i = y_i - \yh_i$ to identify observations that lie
    unusually above or below the regression line.
  \end{definitionbox}

  \medskip

  \textbf{Applications:}
  \begin{itemize}
    \item \textbf{Housing:} the most negative residual is the most ``underpriced'' home.
    \item \textbf{Law schools:} rank by residual after controlling for inputs (LSAT,
      GPA) --- the highest-residual school adds the most ``value.''
    \item \textbf{Sports:} identify over- or under-paid athletes given performance.
  \end{itemize}

\end{frame}

\begin{frame}{Residual Analysis: Housing Example}

  \begin{examplebox}{Housing Prices (HPRICE1)}
    Regress $price$ on $lotsize$, $sqrft$, $bdrms$ ($n = 88$ homes). The most negative
    residual is $\uh_{81} = -120{,}206$: house 81's asking price is about \$120{,}206
    below its predicted price. Possible reasons: unobserved defects, a motivated seller,
    or a genuine bargain.
  \end{examplebox}

  \medskip

  \begin{keybox}
    The residual reflects everything \textbf{not} captured by the model; large residuals
    may signal misspecification or unusual observations.
  \end{keybox}

\end{frame}

\begin{frame}[shrink=8]{Predicting $y$ When $\log(y)$ Is the Dependent Variable}

  If the model is $\log(y) = \beta_0 + \beta_1 x_1 + \cdots + \beta_k x_k + u$,
  OLS predicts $\widehat{\log(y)}$, not $y$ directly.

  \medskip

  \textbf{Naive prediction:} $\tilde{y} = \exp(\widehat{\log(y)})$

  \medskip

  \begin{keybox}
    This \textbf{systematically underestimates} $\E[y \mid \mathbf{x}]$
    because $\E[\exp(u)] > \exp(\E[u]) = 1$ by Jensen's inequality
    whenever $\Var(u) > 0$.
  \end{keybox}

  \medskip

  \begin{definitionbox}{Correction Under Normality (eq.\ 6.40)}
    If $u \mid \mathbf{x} \sim \Normal(0, \sigma^2)$, then
    $\E[\exp(u)] = \exp(\sigma^2/2)$, so:
    \[
      \hat{y} = \exp\!\left(\frac{\hat\sigma^2}{2}\right) \cdot \exp\!\left(\widehat{\log(y)}\right)
    \]
    where $\hat\sigma^2 = \SSR/(n - k - 1)$ is the squared SER.
    The correction factor $\exp(\hat\sigma^2/2) > 1$ inflates the naive prediction.
  \end{definitionbox}

\end{frame}

\begin{frame}[shrink=8]{Predicting $y$ From $\log(y)$: Smearing Estimate}

  \begin{definitionbox}{Smearing Estimate (eq.\ 6.43)}
    Without assuming normality, estimate $\E[\exp(u)]$ from residuals:
    \[
      \hat\alpha_0 = \frac{1}{n} \sum_{i=1}^{n} \exp(\uh_i)
    \]
    Then the adjusted prediction is:
    \[
      \hat{y} = \hat\alpha_0 \cdot \exp\!\left(\widehat{\log(y)}\right)
    \]
    Since $\exp$ is convex, $\hat\alpha_0 > 1$ always.
  \end{definitionbox}

  \medskip

  \begin{keybox}
    \textbf{Which correction to use?}
    \begin{itemize}
      \item Errors approximately normal $\Rightarrow$ use eq.\ 6.40
            ($\exp(\hat\sigma^2/2)$); more efficient when correct
      \item Unsure about normality $\Rightarrow$ use smearing estimate (eq.\ 6.43);
            robust to non-normality
      \item Both are consistent estimators of $\E[y \mid \mathbf{x}]$
    \end{itemize}
  \end{keybox}

\end{frame}

\begin{frame}[shrink=8]{Example 6.7 --- Predicting CEO Salaries}

  \begin{examplebox}{CEO Salary Prediction (CEOSAL2 --- eq.\ 6.45)}
    \[
      \widehat{lsalary} = \underset{(0.257)}{4.504}
        + \underset{(0.039)}{0.163}\, lsales
        + \underset{(0.050)}{0.109}\, lmktval
        + \underset{(0.0053)}{0.0117}\, ceoten
    \]
    \[
      n = 177, \quad \Rsq = 0.318
    \]
  \end{examplebox}

  \medskip

      \textbf{Predict for:} $sales = 5{,}000$ (millions), \\
      $mktval = 10{,}000$ (millions), $ceoten = 10$ years.

      \medskip

      $\widehat{lsalary} = 4.504 + 0.163 \log(5{,}000)$
      $+ 0.109 \log(10{,}000) + 0.0117(10) \approx 7.013$

      \medskip

      Naive: $\exp(7.013) = \$1{,}110{,}983$
  \medskip

      \begin{keybox}
        \textbf{Adjusted predictions:}
        \begin{itemize}
          \item Smearing: $\hat\alpha_0 = 1.136$ \\
                $\hat{y} = 1.136 \times 1{,}110{,}983 = \$1{,}262{,}077$
          \item Normality: $\exp(\hat\sigma^2/2) = 1.117$ \\
                $\hat{y} = 1.117 \times 1{,}110{,}983 = \$1{,}240{,}968$
        \end{itemize}

        \medskip

        Both adjustments add $\approx$ \$130,000--\$150,000
        to the naive prediction. The adjustment matters!
      \end{keybox}

\end{frame}

\begin{frame}[shrink=8]{Comparing Models with Different Dependent Variables}

      \textbf{Problem:} Cannot compare $\Rsq$ from a model
      for $y$ with $\Rsq$ from a model for $\log(y)$.

      \medskip

      \textbf{Solution:} Compute a ``pseudo-$\Rsq$'' for $y$
      from the log model:
      \[
        \tilde{R}^2 = 1 - \frac{\sum_{i=1}^{n} \hat{r}_i^2}
                             {\sum_{i=1}^{n} (y_i - \bar{y})^2}
      \]
      where $\hat{r}_i = y_i - \hat\alpha_0 \exp(\widehat{\log(y)}_i)$
      are residuals for predicting $y$ (not $\log y$).
  \medskip

      \begin{examplebox}{Example 6.8: CEO Salaries}
        From the log model:
        $\Corr(salary_i, \hat{m}_i)^2 = (0.493)^2 = 0.243$

        \medskip

        From the level model (eq.\ 6.49):
        $\Rsq = 0.201$

        \medskip

        The log model explains \textbf{more} variation
        in $salary$ (24.3\% vs.\ 20.1\%) and has more
        interpretable coefficients.
      \end{examplebox}

      \medskip

      \medskip
      \textbf{Conclusion.} The log model wins on both goodness-of-fit and
      interpretability for CEO salary.

\end{frame}

%% ============================================================
\section*{Summary}
%% ============================================================

\begin{frame}[shrink=8]{Chapter 6 Summary: Scaling and Functional Form}

      \begin{keybox}
        \textbf{Data Scaling (\S6-1)}
        \begin{itemize}
          \item Rescaling changes $\hat\beta$ and SEs proportionally
          \item $t$, $F$, $\Rsq$ are \textbf{invariant} to rescaling
          \item Beta coefficients: standardize all variables for
                cross-regressor comparison
        \end{itemize}
      \end{keybox}
  \medskip

      \begin{keybox}
        \textbf{Functional Form (\S6-2)}
        \begin{itemize}
          \item Four log combinations; exact \% change $= \exp(\cdot) - 1$
          \item Quadratics: marginal effect $= \hat\beta_1 + 2\hat\beta_2 x$
          \item Interactions: partial effect depends on other variable
          \item APE: average partial effect over sample
          \item Centering at $\bar{x}$: constant $=$ APE directly
        \end{itemize}
      \end{keybox}

\end{frame}

\begin{frame}[shrink=8]{Chapter 6 Summary: Model Selection and Prediction}

      \begin{keybox}
        \textbf{Adjusted $\Rsq$ (\S6-3)}
        \begin{itemize}
          \item Penalizes extra regressors; can be negative
          \item Valid for nonnested models with \textbf{same} $y$
          \item Cannot compare models with different $y$'s
          \item Beware over-controlling (bad controls)
        \end{itemize}
      \end{keybox}
  \medskip

      \begin{keybox}
        \textbf{Prediction (\S6-4)}
        \begin{itemize}
          \item CI for $\E[y|\mathbf{x}]$: narrow (sampling uncertainty)
          \item PI for new $y^0$: wide (adds $\hat\sigma^2$ for error)
          \item Log models: correct with $\exp(\hat\sigma^2/2)$
                or smearing estimate
          \item Compare $y$ vs.\ $\log(y)$ using pseudo-$\Rsq$
        \end{itemize}
      \end{keybox}

  \medskip

  \begin{highlightbox}
    \textbf{Core theme of Ch.\ 6:} OLS is more flexible than it looks.
    Scaling, functional form, and prediction all build on the same estimator ---
    what changes is \textit{how we set up the model} and \textit{how we read the results}.
  \end{highlightbox}

\end{frame}

%% ============================================================
\section*{Appendix 6A: Bootstrap}
%% ============================================================

\begin{frame}[shrink=8]{The Bootstrap: Motivation}

      \textbf{Problem:} Formulas for standard errors are sometimes
      \begin{itemize}
        \item Hard to derive analytically
        \item Not good approximations to the true sampling variation
        \item Unavailable for complex estimators
      \end{itemize}

      \medskip

      \textbf{Solution:} Use a \textbf{resampling method} --- treat
      the observed data as a population from which we can draw
      new samples.

      \medskip

      \begin{keybox}
        The most common resampling method is the
        \textbf{nonparametric bootstrap}.
        It requires no distributional assumptions.
      \end{keybox}
  \medskip

      \begin{definitionbox}{The Bootstrap Idea}
        \begin{enumerate}
          \item We have an estimator $\hat{\theta}$ computed
                from a sample of size $n$
          \item Treat the $n$ observations as the ``population''
          \item Draw $n$ observations \textbf{with replacement}
                to create a bootstrap sample
          \item Recompute $\hat{\theta}^{(b)}$ on each bootstrap sample
          \item Repeat $m$ times and use the variation in
                $\{\hat{\theta}^{(b)}\}$ to measure uncertainty
        \end{enumerate}
      \end{definitionbox}

\end{frame}

\begin{frame}{Bootstrap Standard Error: The Algorithm}

  \textbf{Algorithm:}
  \begin{enumerate}
    \item list observations $1, \ldots, n$;
    \item draw $n$ integers with replacement $\Rightarrow$ bootstrap sample $b$ (size $n$);
    \item compute the estimate $\hat{\theta}^{(b)}$ on sample $b$;
    \item repeat steps 2--3 a total of $m$ times;
    \item compute $\bar{\theta} = m^{-1}\sum_{b=1}^m \hat{\theta}^{(b)}$, then
      $\text{bse}(\hat{\theta})$ from eq.\ [6.50].
  \end{enumerate}

  \medskip

  \begin{keybox}
    \textbf{Choice of $m$:} typically $m = 1{,}000$ (even $m = 500$ is often reliable).
    Note that $m$ has nothing to do with the sample size $n$.
  \end{keybox}

\end{frame}

\begin{frame}{Bootstrap Standard Error: The Formula}

  \begin{definitionbox}{Bootstrap Standard Error (eq.\ 6.50)}
    \[
      \text{bse}(\hat{\theta}) = \left[
        (m-1)^{-1} \sum_{b=1}^{m}
        \bigl(\hat{\theta}^{(b)} - \bar{\theta}\bigr)^2
      \right]^{1/2},
    \]
    where $\bar{\theta} = m^{-1}\sum_{b=1}^m \hat{\theta}^{(b)}$ is the average of the
    bootstrap estimates.
  \end{definitionbox}

  \medskip

  \begin{highlightbox}
    $\text{bse}(\hat{\theta})$ is just the \textbf{sample standard deviation} of the $m$
    bootstrap estimates --- a consistent estimator of $\text{se}(\hat{\theta})$.
  \end{highlightbox}

\end{frame}

\begin{frame}{Bootstrap in Practice}

  \begin{keybox}
    \textbf{When to use the bootstrap:}
    \begin{itemize}
      \item the asymptotic SE formula is unavailable (e.g.\ nonlinear functions of
        $\hat\beta$);
      \item small samples where asymptotics may be poor;
      \item checking whether analytical SEs are reliable.
    \end{itemize}
  \end{keybox}

  \medskip

  \begin{keybox}
    \textbf{Going further.} Bootstrap samples can also compute \textit{p}-values for $t$
    and $F$ statistics, or construct confidence intervals directly --- often more
    accurate than putting $\text{bse}(\hat{\theta})$ into a $t$ statistic. See
    Horowitz (2001).
  \end{keybox}

\end{frame}

%% ============================================================
\end{document}
